\IfFileExists{stacks-project.cls}{%
\documentclass{stacks-project}
}{%
\documentclass{amsart}
}
\usepackage{amssymb}

% For dealing with references we use the comment environment
\usepackage{verbatim}
\usepackage{fancyvrb}
\usepackage{pmboxdraw}
\newenvironment{reference}{\comment}{\endcomment}
%\newenvironment{reference}{}{}
\newenvironment{slogan}{\comment}{\endcomment}
\newenvironment{history}{\comment}{\endcomment}

% For source figures retained as assets
\usepackage{graphicx}

% For commutative diagrams we use Xy-pic
\usepackage[all]{xy}

% We use 2cell for 2-commutative diagrams.
\xyoption{2cell}
\UseAllTwocells

% We use multicol for the list of chapters between chapters
\usepackage{multicol}

% This is generally recommended for better output
\usepackage{lmodern}
\usepackage[T1]{fontenc}
\usepackage{textalpha}

% For cross-file-references
% We prefer the snapshot tmp/xref/#2.aux,
% which latexmkrc replaces atomically, so
% that documents compiled simultaneously
% never read a half-written .aux file.
% The last argument keeps links to #2.pdf.
\usepackage{xr-hyper}
\newcommand{\maybeexternaldocument}[2]{%
  \IfFileExists{#2.aux}{%
    \externaldocument[#1]{#2}%
  }{%
    \PackageWarning{stack}{Missing #2.aux; external references with prefix #1 unavailable}%
  }%
}
\let\externaldocumentorig\externaldocument
\renewcommand{\externaldocument}[2][]{%
  \IfFileExists{tmp/xref/#2.aux}{%
    \externaldocumentorig[#1]{tmp/xref/#2}[#2.pdf]%
  }{%
  \IfFileExists{#2.aux}{%
    \externaldocumentorig[#1]{#2}%
  }{%
    \PackageWarning{stack}{Missing #2.aux; external references with prefix #1 unavailable}%
  }%
  }%
}

% Package for hypertext links:
\usepackage{hyperref}

% For any local file, say "hello.tex" you want to link to please
% use \externaldocument[hello-]{hello}
\externaldocument[introduction-]{introduction}
\externaldocument[conventions-]{conventions}
\externaldocument[sets-]{sets}
\externaldocument[categories-]{categories}
\externaldocument[topology-]{topology}
\externaldocument[sheaves-]{sheaves}
\externaldocument[sites-]{sites}
\externaldocument[stacks-]{stacks}
\externaldocument[fields-]{fields}
\externaldocument[algebra-]{algebra}
\externaldocument[brauer-]{brauer}
\externaldocument[homology-]{homology}
\externaldocument[derived-]{derived}
\externaldocument[simplicial-]{simplicial}
\externaldocument[more-algebra-]{more-algebra}
\externaldocument[smoothing-]{smoothing}
\externaldocument[modules-]{modules}
\externaldocument[sites-modules-]{sites-modules}
\externaldocument[injectives-]{injectives}
\externaldocument[cohomology-]{cohomology}
\externaldocument[sites-cohomology-]{sites-cohomology}
\externaldocument[dga-]{dga}
\externaldocument[dpa-]{dpa}
\externaldocument[sdga-]{sdga}
\externaldocument[hypercovering-]{hypercovering}
\externaldocument[schemes-]{schemes}
\externaldocument[constructions-]{constructions}
\externaldocument[properties-]{properties}
\externaldocument[morphisms-]{morphisms}
\externaldocument[coherent-]{coherent}
\externaldocument[divisors-]{divisors}
\externaldocument[limits-]{limits}
\externaldocument[varieties-]{varieties}
\externaldocument[topologies-]{topologies}
\externaldocument[descent-]{descent}
\externaldocument[perfect-]{perfect}
\externaldocument[more-morphisms-]{more-morphisms}
\externaldocument[flat-]{flat}
\externaldocument[groupoids-]{groupoids}
\externaldocument[more-groupoids-]{more-groupoids}
\externaldocument[etale-]{etale}
\externaldocument[chow-]{chow}
\externaldocument[intersection-]{intersection}
\externaldocument[pic-]{pic}
\externaldocument[weil-]{weil}
\externaldocument[adequate-]{adequate}
\externaldocument[dualizing-]{dualizing}
\externaldocument[duality-]{duality}
\externaldocument[discriminant-]{discriminant}
\externaldocument[derham-]{derham}
\externaldocument[local-cohomology-]{local-cohomology}
\externaldocument[algebraization-]{algebraization}
\externaldocument[curves-]{curves}
\externaldocument[resolve-]{resolve}
\externaldocument[models-]{models}
\externaldocument[functors-]{functors}
\externaldocument[equiv-]{equiv}
\externaldocument[pione-]{pione}
\externaldocument[etale-cohomology-]{etale-cohomology}
\externaldocument[proetale-]{proetale}
\externaldocument[relative-cycles-]{relative-cycles}
\externaldocument[more-etale-]{more-etale}
\externaldocument[trace-]{trace}
\externaldocument[crystalline-]{crystalline}
\externaldocument[spaces-]{spaces}
\externaldocument[spaces-properties-]{spaces-properties}
\externaldocument[spaces-morphisms-]{spaces-morphisms}
\externaldocument[decent-spaces-]{decent-spaces}
\externaldocument[spaces-cohomology-]{spaces-cohomology}
\externaldocument[spaces-limits-]{spaces-limits}
\externaldocument[spaces-divisors-]{spaces-divisors}
\externaldocument[spaces-over-fields-]{spaces-over-fields}
\externaldocument[spaces-topologies-]{spaces-topologies}
\externaldocument[spaces-descent-]{spaces-descent}
\externaldocument[spaces-perfect-]{spaces-perfect}
\externaldocument[spaces-more-morphisms-]{spaces-more-morphisms}
\externaldocument[spaces-flat-]{spaces-flat}
\externaldocument[spaces-groupoids-]{spaces-groupoids}
\externaldocument[spaces-more-groupoids-]{spaces-more-groupoids}
\externaldocument[bootstrap-]{bootstrap}
\externaldocument[spaces-pushouts-]{spaces-pushouts}
\externaldocument[spaces-chow-]{spaces-chow}
\externaldocument[groupoids-quotients-]{groupoids-quotients}
\externaldocument[spaces-more-cohomology-]{spaces-more-cohomology}
\externaldocument[spaces-simplicial-]{spaces-simplicial}
\externaldocument[spaces-duality-]{spaces-duality}
\externaldocument[formal-spaces-]{formal-spaces}
\externaldocument[restricted-]{restricted}
\externaldocument[spaces-resolve-]{spaces-resolve}
\externaldocument[formal-defos-]{formal-defos}
\externaldocument[defos-]{defos}
\externaldocument[cotangent-]{cotangent}
\externaldocument[examples-defos-]{examples-defos}
\externaldocument[algebraic-]{algebraic}
\externaldocument[examples-stacks-]{examples-stacks}
\externaldocument[stacks-sheaves-]{stacks-sheaves}
\externaldocument[criteria-]{criteria}
\externaldocument[artin-]{artin}
\externaldocument[quot-]{quot}
\externaldocument[stacks-properties-]{stacks-properties}
\externaldocument[stacks-morphisms-]{stacks-morphisms}
\externaldocument[stacks-limits-]{stacks-limits}
\externaldocument[stacks-cohomology-]{stacks-cohomology}
\externaldocument[stacks-perfect-]{stacks-perfect}
\externaldocument[stacks-introduction-]{stacks-introduction}
\externaldocument[stacks-more-morphisms-]{stacks-more-morphisms}
\externaldocument[stacks-geometry-]{stacks-geometry}
\externaldocument[moduli-]{moduli}
\externaldocument[moduli-curves-]{moduli-curves}
\externaldocument[examples-]{examples}
\externaldocument[exercises-]{exercises}
\externaldocument[guide-]{guide}
\externaldocument[desirables-]{desirables}
\externaldocument[coding-]{coding}
\externaldocument[obsolete-]{obsolete}
\externaldocument[fdl-]{fdl}
\externaldocument[index-]{index}

% Heap Project documents
\externaldocument[k3-]{k3}
\externaldocument[torelli-]{torelli}
\externaldocument[linear-algebra-]{linear-algebra}
\externaldocument[groups-]{groups}
\externaldocument[complex-analysis-]{complex-analysis}
\externaldocument[functional-analysis-]{functional-analysis}
\externaldocument[categories-]{categories}
\externaldocument[commutative-algebra-]{commutative-algebra}
\externaldocument[ringed-spaces-]{ringed-spaces}
\externaldocument[homological-algebra-]{homological-algebra}
\externaldocument[lie-algebras-]{lie-algebras}
\externaldocument[representation-theory-]{representation-theory}
\externaldocument[smooth-manifolds-]{smooth-manifolds}
\externaldocument[differential-topology-]{differential-topology}
\externaldocument[algebraic-topology-]{algebraic-topology}
\externaldocument[differential-geometry-]{differential-geometry}
\externaldocument[algebraic-geometry-]{algebraic-geometry}
\externaldocument[symplectic-geometry-]{symplectic-geometry}
\externaldocument[complex-geometry-]{complex-geometry}
\externaldocument[hodge-theory-]{hodge-theory}
\externaldocument[%
variations-of-hodge-structures-]{%
variations-of-hodge-structures}
\externaldocument[%
homological-mirror-symmetry-]{%
homological-mirror-symmetry}
\maybeexternaldocument{deformations-}{deformations}
\externaldocument[vertex-algebras-]{vertex-algebras}
\externaldocument[springer-]{springer}
\externaldocument[infty-categories-]{infty-categories}
\externaldocument[seminars-differential-geometry-]{%
seminars-differential-geometry}
\externaldocument[seminars-complex-geometry-]{%
seminars-complex-geometry}
\externaldocument[seminars-algebraic-geometry-]{%
seminars-algebraic-geometry}
\externaldocument[seminars-analysis-pde-]{%
seminars-analysis-pde}
\externaldocument[surfaces-]{surfaces}
\externaldocument[birational-maps-commutative-algebra-]{birational-maps-commutative-algebra}
\externaldocument[cremona-transformations-]{cremona-transformations}
\externaldocument[derived-categories-of-sheaves-]{derived-categories-of-sheaves}
\externaldocument[geometric-invariant-theory-]{geometric-invariant-theory}
\externaldocument[geometric-stability-conditions-and-group-actions-]{geometric-stability-conditions-and-group-actions}
\externaldocument[moduli-spaces-of-sheaves-]{moduli-spaces-of-sheaves}
\externaldocument[bridgeland-stability-general-theory-]{bridgeland-stability-general-theory}
\externaldocument[hilbert-schemes-]{hilbert-schemes}
\externaldocument[noncommutative-abelian-surfaces-and-kummer-type-hyperkahler-varieties-]{noncommutative-abelian-surfaces-and-kummer-type-hyperkahler-varieties}
\externaldocument[mumford-tate-groups-in-hodge-theory-]{mumford-tate-groups-in-hodge-theory}
\externaldocument[hodge-birational-atoms-2026-]{hodge-birational-atoms-2026}
\externaldocument[xxii-egd-2026-]{%
xxii-egd-2026}
\externaldocument[pde-]{pde}
\externaldocument[probability-]{probability}

\externaldocument[linguistics-]{linguistics}

\externaldocument[%
16th-alga-meeting-2026-]{%
16th-alga-meeting-2026}

\externaldocument[%
iv-jornadas-lefschetz-2026-]{%
iv-jornadas-lefschetz-2026}

\externaldocument[reading-the-masters-]{%
reading-the-masters}

\externaldocument[real-analysis-]{%
real-analysis}

\externaldocument[optimal-transport-]{%
optimal-transport}

\externaldocument[history-]{history}

\externaldocument[seiberg-witten-]{seiberg-witten}

\externaldocument[history-philosophy-biology-]{history-philosophy-biology}

\externaldocument[classical-mechanics-]{classical-mechanics}
\externaldocument[electromagnetism-]{electromagnetism}
\externaldocument[waves-and-optics-]{waves-and-optics}
\externaldocument[wif-]{waves-in-fluids-iv}
\externaldocument[thermal-physics-]{thermal-physics}
\externaldocument[quantum-mechanics-]{quantum-mechanics}
\externaldocument[special-relativity-]{special-relativity}
\externaldocument[physics-]{physics}

\externaldocument[chemistry-]{chemistry}

\externaldocument[molecular-cell-biology-]{molecular-cell-biology}

\externaldocument[microbiology-]{microbiology}

\externaldocument[protistology-]{protistology}

\externaldocument[mycology-]{mycology}

\externaldocument[botany-]{botany}

\externaldocument[zoology-]{zoology}

\externaldocument[animal-histology-]{animal-histology}

\externaldocument[animal-physiology-]{animal-physiology}

\externaldocument[systematics-biogeography-]{systematics-biogeography}

\externaldocument[ecology-natural-resources-]{ecology-natural-resources}

\externaldocument[earth-sciences-]{earth-sciences}

\externaldocument[biostatistics-]{biostatistics}

\externaldocument[mathematical-systems-biology-]{mathematical-systems-biology}

\externaldocument[cimpa-el-salvador-2024-]{cimpa-el-salvador-2024}

\externaldocument[anthropology-]{anthropology}

\externaldocument[yang-mills-theory-]{yang-mills-theory}

\externaldocument[seminars-algebra-]{seminars-algebra}

\externaldocument[seminars-physics-]{seminars-physics}

\externaldocument[seminars-computing-]{seminars-computing}

\externaldocument[oktobermat-2026-]{oktobermat-2026}

\externaldocument[saet-2026-]{saet-2026}

\externaldocument[seminars-symplectic-geometry-]{seminars-symplectic-geometry}

\externaldocument[seminars-dynamics-]{seminars-dynamics}

\externaldocument[oktobermat-2024-]{oktobermat-2024}

\externaldocument[quantum-groups-]{quantum-groups}

\externaldocument[legal-2027-]{legal-2027}

\externaldocument[derived-categories-]{derived-categories}

% Theorem environments.
%
\theoremstyle{plain}
\newtheorem{theorem}[subsection]{Theorem}
\newtheorem{proposition}[subsection]{Proposition}
\newtheorem{lemma}[subsection]{Lemma}
\newtheorem{corollary}[subsection]{Corollary}

\theoremstyle{definition}
\newtheorem{definition}[subsection]{Definition}
\newtheorem{example}[subsection]{Example}
\newtheorem{exercise}[subsection]{Exercise}
\newtheorem{situation}[subsection]{Situation}

\theoremstyle{remark}
\newtheorem{remark}[subsection]{Remark}
\newtheorem{remarks}[subsection]{Remarks}

\numberwithin{equation}{subsection}

% Macros
%
\def\lim{\mathop{\mathrm{lim}}\nolimits}
\def\colim{\mathop{\mathrm{colim}}\nolimits}
\def\Spec{\mathop{\mathrm{Spec}}}
\def\Hom{\mathop{\mathrm{Hom}}\nolimits}
\def\Ext{\mathop{\mathrm{Ext}}\nolimits}
\def\SheafHom{\mathop{\mathcal{H}\!\mathit{om}}\nolimits}
\def\SheafExt{\mathop{\mathcal{E}\!\mathit{xt}}\nolimits}
\def\Sch{\mathit{Sch}}
\def\Mor{\mathop{\mathrm{Mor}}\nolimits}
\def\Ob{\mathop{\mathrm{Ob}}\nolimits}
\def\Sh{\mathop{\mathit{Sh}}\nolimits}
\def\NL{\mathop{N\!L}\nolimits}
\def\CH{\mathop{\mathrm{CH}}\nolimits}
\def\proetale{{pro\text{-}\acute{e}tale}}
\def\etale{{\acute{e}tale}}
\def\QCoh{\mathit{QCoh}}
\def\Ker{\mathop{\mathrm{Ker}}}
\def\Im{\mathop{\mathrm{Im}}}
\def\Coker{\mathop{\mathrm{Coker}}}
\def\Coim{\mathop{\mathrm{Coim}}}

% Boxtimes
%
\DeclareMathSymbol{\boxtimes}{\mathbin}{AMSa}{"02}

%
% Macros for moduli stacks/spaces
%
\def\QCohstack{\mathcal{QC}\!\mathit{oh}}
\def\Cohstack{\mathcal{C}\!\mathit{oh}}
\def\Spacesstack{\mathcal{S}\!\mathit{paces}}
\def\Quotfunctor{\mathrm{Quot}}
\def\Hilbfunctor{\mathrm{Hilb}}
\def\Curvesstack{\mathcal{C}\!\mathit{urves}}
\def\Polarizedstack{\mathcal{P}\!\mathit{olarized}}
\def\Complexesstack{\mathcal{C}\!\mathit{omplexes}}
% \Pic is the operator that assigns to X its picard group, usage \Pic(X)
% \Picardstack_{X/B} denotes the Picard stack of X over B
% \Picardfunctor_{X/B} denotes the Picard functor of X over B
\def\Pic{\mathop{\mathrm{Pic}}\nolimits}
\def\Picardstack{\mathcal{P}\!\mathit{ic}}
\def\Picardfunctor{\mathrm{Pic}}
\def\Deformationcategory{\mathcal{D}\!\mathit{ef}}


\begin{document}

\title{Philosophy}
\maketitle

\phantomsection
\label{section-phantom}

\tableofcontents

\section{Courses and sources}
\label{section-courses-and-sources}

These are my texts from the philosophy
courses I took at the Facultad de
Filosofía y Letras of UNAM in 2017 and
2018, while studying mathematics. They are
kept in Spanish. I have edited them lightly
for publication: informal asides and
personal remarks were removed, and the
arguments are left as they were, including
their slips. Greek words that I had typed in
Greek script are transliterated here.

The texts come from four courses and one
reading.
\begin{enumerate}
\item A reading of Ricoeur's \emph{Tiempo y
narración} in the first half of 2017. I
have a single note from it
(Section~\ref{section-ricoeur-tiempo}).
The course was taught by Greta Rivara; its
name is not recorded in my files.
\item Historia de la Filosofía I, second
half of 2017: the Presocratics and Plato
(Sections~\ref{section-jenofanes} and
\ref{section-teeteto}).
\item Textos Filosóficos I, second half of
2017: Plato's \emph{Republic}
(Sections~\ref{section-republica-vi-vii} and
\ref{section-republica-democracia}).
\item Historia de la Filosofía II, semester
2018-2, taught by María Teresa Padilla
Longoria: Aristotle and the Hellenistic
schools
(Sections~\ref{section-notas-aristoteles},
\ref{section-termino-medio} and
\ref{section-sabio-cinico}).
\item Filosofía del Lenguaje, semester
2018-2, taught by Erika Rebeca Lindig
Cisneros
(Sections~\ref{section-farmacia-de-platon}
and \ref{section-derrida-juego}). The
readings of that course are described in
Section~\ref{section-language}.
\end{enumerate}

\section{Ricoeur: el tiempo lineal es R}
\label{section-ricoeur-tiempo}

A note of June 2017 on the first chapter of
the third volume of Ricoeur's \emph{Tiempo
y narración}, which sets Aristotle's time
of the world against Augustine's time of
the soul. The note compares Aristotle's
continuum with the real line.

Dice Ricoeur: a la concepción fenomenológica
del tiempo debemos pagar el precio de la
creciente aporía.

Aristóteles: el movimiento implica tiempo. El
movimiento está en las cosas que cambian,
mientras el tiempo está en todas partes y
emerge del movimiento.

Aristóteles: hay tres cosas: magnitud,
movimiento y tiempo; lo que cambia, el
cambio, el antes y el después de emergen de
todo esto. Esta relación de orden (antes y
después) ``está en el mundo mucho antes de
estar en el alma''.

Ricoeur: la continuidad es la propiedad de
algo de ser fragmentado en partes
infinitamente. Así, el tiempo es continuo.

``La definición aristotélica del tiempo y los
números reales.'' Tomo III, p. 649

\begin{quote}
La continuidad es la propiedad de algo de ser
fragmentado en partes infinitamente. La
continuidad de las cosas que cambian
(magnitud), del cambio (movimiento) y del
tiempo (el antes y el después).

El número es la medida que se le asigna al
intervalo entre el antes y el después.

¿La percepción humana es indispensable para
esta construcción del tiempo? Según
Aristóteles, no.
\end{quote}

El número es la medida que se le asigna al
intervalo entre el antes y el después; eso es
el tiempo. En este sentido, de acuerdo a la
tesis de Aristóteles de que el tiempo emerge
del movimiento, lo que nos importa no es
tanto el número (el tiempo) sino la propiedad
de ser numerable (del movimiento).

Y luego tira ésta, brutal:

``Para ser pensable, el instante de
Aristóteles solamente requiere de un corte
realizado por el espíritu en la continuidad
del movimiento, en tanto éste es numerable''

Así, cada punto en la recta real es digno de
ser el presente; cada instante puede estar
sujeto a este corte del espíritu. Pero ``el
presente agustinano es cualquier instante
designado por el hablante como el ``ahora de
su encunciación.''

El presente fenomenológico no es cualquier
punto de la recta real, sino que es el
momento en el que yo soy.

Una observación matemática: una de las
propiedades más místicas de
los números reales es la de continuidad. Con
el teorema de los intervalos anidados, se
demuestra que este conjunto ordenado
linealmente ``no tiene hoyos'', siendo así
que por más que esté constituido por puntos
(o instantes) individuales y distintos, se
nos figura un todo ``bien construido'', casi
real.

Ricoeur define el continuo por su propiedad
de ser divisible infinitamente. Para
nosotros, es importante sacudir la ambigüedad
en esta definición. El ejemplo más claro es
el conjunto de los números racionales, donde
podemos ``dividir infinitamente'' un
intervalo $[a,b]$ sin jamás llegar al
conjunto vacío; pero, sabemos que Q ``sí
tiene hoyitos''.

Todo estriba en esa pregunta: ¿qué quiere
decir que un conjunto ordenado tenga hoyitos?
En teoría de conjuntos, llamamos ``la
propiedad de la mínima cota superior'' a la
implicación de que todo subconjunto no vacío
acotado superiormente tenga supremo. Spivak
propone esta propiedad matemática como un
axioma, del cual se desprende la completitud
(o continuidad) de R.

\medskip\noindent\emph{Aristóteles y la
derivada:}

Magnitud, movimiento y tiempo. Lo que cambia,
el cambio y el antes y el después. La
magnitud es una función. El movimiento
(cambio) es la diferencia función evaluada en
dos puntos distintos, $f(a)-f(b)$. El tiempo
es $a-b$ (el número).

La derivada es el cambio con respecto al
tiempo en un intervalo infinitamente pequeño.

Una pieza linda: ``El instante acompaña al
móvil como el tiempo al movimiento''.

Y luego otra: ``Con otras palabras, el
instante sería en un sentido el presente, en
otro sentido un punto del tiempo, el presente
siempre el mismo recorriendo puntos de tiempo
continuamente diferentes''.

\section{La crítica religiosa de Jenófanes}
\label{section-jenofanes}

An answer written in October 2017 for
Historia de la Filosofía I to the question
``¿En qué consiste la crítica religiosa de
Jenófanes?''

Me parece que los poemas de Jenófanes
representan un cambio en la noción de dios en
la antigua Grecia. Los poemas de Homero y de
Hesíodo influyeron de tal modo al pueblo
griego que su cosmogonía y sus nociones de
qué eran los dioses se tomaban con mucha
seriedad, y la crítica religiosa de Jenófanes
viene a contrariar esta confianza en los
poetas antiguos.

Yo pienso que hay que preguntarnos dos cosas:
¿por qué fueron tan influyentes los poemas
épicos?, y ¿en qué estriba la necesidad de
Jenófanes de criticarlos?

Homero y Hesíodo se nos figuran como los
primeros poetas de la antigua Grecia. Sus
poemas traen a nosotros no sólo una muestra
de altas capacidades literarias sino un
horizonte filosófico que debemos valorar
mucho: por un lado, nos dan una cosmogonía
muy completa, y por otro, nos hablan de
conceptos abstractos, naturales y humanos,
como el caos y la belleza, representados
alegóricamente a través de ``dioses''.
Durante la lectura de estos poemas, parece
que se nos enredan hasta lo inseparable estas
nociones abstractas con los personajes
antropomorfos que las representan.

Me parece que Jenófanes restringe la noción
de dios a lo que podemos llamar un ente
supremo que se nos figura como límite o base
de todo lo existente: ``Uno solo es dios,
entre hombres y dioses el más grande, ni un
cuerpo parejo a los mortales, ni en
entendimiento.''; ``Todo él ve, todo él
entiende, todo él oye,''. Uso la palabra
``restringe'' porque me parece que ya Zeus
había jugado ese papel, como muestro a
continuación.

Por ejemplo, tenemos en la Teogonía, en el
verso 613 respecto del mito de Prometeo:
``Así, no es posible engañar ni eludir la
mente de Zeus'', y, repetidamente en esta
parte del poema: ``Zeus, que conoce
inmortales consejos''\footnote{Teogonía,
Bibliotheca Scriptorum Graecorum et Romanorum
Mexicana}. Vemos en estas citas una especie
de panlogismo en la figura de Zeus. En un
plano más general, otra defensa de este
argumento es la supremacía de Zeus sobre los
demás dioses, que representa, más allá del
mito, una diferencia ontológica fundamental
entre las distintas figuras divinas. Y
finalmente, en el plano lingüístico, no
podemos ignorar el hecho de que la palabra
``Zeus'' proviene de la indoeuropea
``Dieus'', previa a la
mitología griega pero sí asociada a un
principio ``supremo''.

Es necesario para Jenófanes librar a este
ente supremo y perfecto de su representación
literaria antropomorfa capaz de cualquier
barbaridad moral, para conseguir dos cosas:
primero, evitar que la lectura de los
estimadísimos poemas antiguos fomentara la
realización de actos poco virtuosos al
mostrar cómo los mismos dioses los cometían,
y luego, para dar un nuevo paso en la
delimitación de aquel concepto de ente
supremo por sí mismo, es decir, sin servirse
de ninguna figura literaria y más bien
encapsularlo en la palabra ``dios''.

En suma, pienso que las condiciones
lingüísticas y epistemológicas que rodearon a
Homero y a Hesíodo los llevaron a discutir lo
que más tarde sería para Jenófanes ``dios''
en un contexto mítico y literario que
difícilmente lo hacía ver distinto del
personaje de una fábula que sufre tantas
pasiones. La crítica religiosa de Jenófanes
consiste en la separación de estos dos
objetos con una finalidad moral, y otra de
avanzar en la definición de un ente supremo
en tanto supremo, cosa que no vemos tal cual
en los poemas antiguos y que por ello
representa un gran paso en el camino de la
filosofía griega.

\section{Preguntas sobre el Teeteto}
\label{section-teeteto}

Five questions on Plato's
\emph{Theaetetus}, with my answers, for
Historia de la Filosofía I (2017).

\medskip\noindent\emph{1. ¿Qué consecuencias
epistemológicas se siguen de la crítica
platónica al homo mensura (hombre como
medida) de Protágoras?}

La crítica platónica al homo mensura lo
invalida por completo, y en ese sentido, se
desprende de ella que la verdad es ajena a la
percepción humana: el hombre no es la medida
de todas las cosas. Así, nos vemos obligados
a suponer ya sea que nada es cognoscible o
que existe otra ``medida'' para generar
conocimiento. Creo que la teoría platónica de
las ideas responde a este problema: al menos
en la Républica, Platón argumenta que el
conocimiento verdadero (que por cierto existe
y él llama nóesis) consiste en la capacidad
de aprehender las ideas de las cosas, o las
cosas en sí mismas. En este sentido, podemos
decir que la medida platónica de las cosas
está en el mundo de las ideas, y como éstas
son entes del todo ajenos a la percepción
humana, es una propuesta distinta a la de
Protágoras.

\medskip\noindent\emph{2. ¿Por qué si se
afirma que la propia percepción es el
criterio exclusivo de verdad, la falsedad
sería imposible?}

Si la propia percepción es el criterio
exclusivo de verdad, cualquier cosa que
consideremos falsa lo es. Pero basta la
existencia de una segunda persona que
considere aquello para nosotros falso,
verdadero, para convertirlo en verdadero. En
este sentido, todo juicio o hecho es
posiblemente verdadero y falso. Aunque el
principio de no contradicción no debe ser
asumido como una máxima irrefutable, ya que
existen cosas que razonablemente se nos
figuran ciertas y falsas a un tiempo, también
existen otras que no tanto así. Para mí, va
más allá de toda opinión humana que $1+1=2$,
dado que estos símbolos representen lo
convencional. En este sentido, existen casos
donde la falsedad implica la falta de verdad,
cosa imposible si la percepción es el
criterio exclusivo de verdad.

\medskip\noindent\emph{3. ¿Qué argumentos da
Sócrates para refutar la tesis protagórica y
asegurar que no es lo mismo conocer que
percibir?}

Sócrates presenta al menos tres intentos para
resolver este problema. Primero, argumenta
que si conocer es lo mismo que percibir,
entonces cualquier sujeto sensible, como un
cerdo, posee conocimiento, y en ese sentido
``el cerdo es medida de todas las cosas''
(161c)\footnote{Platón; Marcelo Boeri (2006).
Teeteto. Editorial Losada}, y nadie supera a
nadie en sabiduría. Sin embargo, esto parece
ser algo así como un juego para Sócrates,
pues rápidamente argumenta él mismo, en el
nombre de Protágoras, que ese razonamiento se
vale más de la probabilidad (eikós) que del
logos (163a). La verdad, aquí no entiendo a
qué se refiere con eso de la probabilidad; y
casi tan pronto como menciona esa palabra,
abandona esa tesis en busca una mejor.

El siguiente intento se basa en el sentido de
la vista y en el recuerdo. Entiendo el
argumento por reducción al absurdo como
sigue: si saber es percibir, entonces ver es
saber, puesto que ver es percibir; así, si
alguien ve algo, entonces lo conoce. Más
cuando ya no lo ve, entonces no le conoce;
esto porque si conocer y percibir son
equivalentes, entonces no percibir es no
conocer (la equivalencia entre percibir y
conocer implica la equivalencia entre no
percibir y no conocer). Luego Sócrates
pregunta a Teeteto si el recordar algo que ya
se sabía de antemano implica saberlo en el
instante del recuerdo, y el joven lo concede.
La contradicción está en que recordar algo
que se ha visto implicaría su conocimiento,
pero al mismo su desconocimiento por ya no
verlo más.

Luego Sócrates hace una defensa a Protágoras
en contra de aquel argumento, donde primero
le reprocha el mal uso de las palabras para
confundir a sus interlocutores, y luego
defiende la compatibilidad de su tesis con la
de la existencia de alguien más sabio que
otro asegurando que la sabiduría consiste en
la capacidad de conseguir cosas buenas: ``lo
que parece bueno y justo a cada ciudad es tal
para ella mientras se forma ese juicio; y el
sabio hace que el bien, y no el mal, sea y
parezca tal a cada ciudadano''
(167c\footnote{Esta cita la saqué de la
edición de la biblioteca Vasconcelos; en la
de Losada es así: ``lo que a cada ciudad le
parece justo y honorable, es así para cada
ciudad, por el tiempo que así lo juzgue. Pero
el sabio es el que hace que cada cosa sea y
parezca benéfica para ellas en lugar de
dañina''.}).

Aunque Sócrates no basa su respuesta en torno
de la noción de bien en la que se basa el
supuesto argumento de Protágoras (cosa que me
parecía más natural), vuelve a hacer una
prueba por reducción al absurdo que es
bastante clara: si percibir es conocer,
entonces lo que parece a cada quién es
exactamente como le parece, y en particular,
la opinión de que percibir no es conocer
sería tan válida como la de que percibir es
conocer (dado que alguien sostenga esta
opinión, cosa que asegura Sócrates). Así, la
tesis Protagórica es verdadera y falsa a un
tiempo, y en particular falsa.

\medskip\noindent\emph{4. ¿Cuáles son las
diferencias metodológicas o procedimentales
que señala Sócrates para poder distinguir al
filósofo del sofista?}

Sócrates muestra a Teeteto cómo un sofista
argumentaría la posibilidad de no saber que
se sabe (165b-165e): si ver es saber, cuando
tapamos un ojo y otro no, no vemos con uno y
vemos con otro, y en este sentido no vemos
que vemos, es decir, no sabemos que sabemos.
Este argumento me parece del todo ridículo, y
precisamente creo que es lo que Sócrates
trata de señalar: el sofista, mediante
confusos argumentos, trata de enredar a su
interlocutor en un discurso que no atiende a
lo esencial del problema con el fin de
aprisionarlo lógicamente para después
reclamar una paga. Un filósofo, en cambio,
cooperaría con su compañero para evitar caer
en razonamientos mal planteados (con el
supuesto de que ambos participantes tienen
posibilidad de errar) y con la finalidad de
alcanzar una solución verdadera y sin esperar
más que eso del diálogo.

\medskip\noindent\emph{5. ¿Qué consecuencias
se siguen de la afirmación de Protágoras de
que no hay realidad alguna que tenga
naturaleza propia?}

La tesis de que no haya realidad alguna que
tenga naturaleza propia implica un
relativismo abrumador, además de la
contradicción inmanente que ya fue explicada.
Creo que uno de los puntos críticos de esto
es que la noción de bien sería también
relativa, y así, lo que para uno es malo para
otro puede ser bueno. Esto resulta, desde mi
punto de vista, en una catástrofe social.
Aunque Protágoras asegura que una ciudad
puede convenir en las nociones de justicia e
injusticia dentro de ella misma, creo que el
hecho de que las cosas no tengan naturaleza
propia está en conflicto con esa posibilidad:
¿qué nos asegura que en una sociedad donde
opiniones contrarias son verdaderas se pueda
llegar al consenso respecto lo que es justo y
lo que no? En mi opinión, la tesis de
Protágoras es insostenible.

\section{La República: libros VI y VII}
\label{section-republica-vi-vii}

A summary of Books VI and VII of the
\emph{Republic}, with my own questions at
the end, written in November 2017.

\medskip\noindent\emph{Libro VI}

El comentario final de Sócrates en el libro V
de la República hace referencia a una
cualidad fundamental de los filósofos: ``los
que se adhieren en todo al ser en sí, y no
filódoxos, amigos de la opinión'' (484a). Y
aunque con esto los interlocutores parecen
quedar satisfechos con la definición de
filósofo hasta entonces construida (``Así,
Glaucón, proseguí, se nos ha mostrado al fin
{[}\ldots{}] quiénes son filósofos'' [484a]),
parece ser que el tema predominante al inicio
del libro VI no es muy distinto: se seguirá
hablando del filósofo, salvo que ahora se
llegará a él a partir de la pregunta:
``[\ldots{}] quiénes habrán de ser los guías
de la ciudad?'' (484a).

Muy pronto en el diálogo, como es costumbre,
vemos surgir una y otra virtud como
características esenciales de aquel que
debería regir la ciudad, y tampoco nos
sorprende el hecho de que estas
características coincidan con lo que
entendemos en el diálogo por filósofo:

\begin{quote}
``[\ldots{}] si no es por naturaleza
memorioso, fácil en el aprender, con nobleza
moral y buena gracia, y además amigo y
allegado de la verdad, de la justicia, de la
fortaleza y de la templanza?

Ni el propio Momo, dijo [Glaucón], podría
censurar a una persona semejante.

Pues bien, proseguí, ¿no sería a tales
hombres, que han llegado así a su plena sazón
por la educación y por su edad, y sólo a
ellos, a quienes confiarías la ciudad?''
(487a)
\end{quote}

Con lo que concluye esta primera
argumentación, dando pie a una queja de
Adimanto: ``podría uno contestarte que
{[}\ldots{}] [los filósofos] resultan en su
mayoría unos tipos del todo extraños, por no
decir del todo perversos'' (487d). Aunque
Adimanto no se atreve a apropiarse ni a
deslindarse de esta postura, Sócrates no duda
en hacerla suya, para después explicar a
través de una analogía que (a grandes rasgos)
en un mundo de ciegos, los videntes son, en
efecto, unos tipos del todo extraños.

De manera similar, Sócrates argumenta que la
inutilidad de un filósofo, que sí existe, lo
hace a cuenta de que aquellos a su alrededor
no saben aprovecharle. Luego, Sócrates se da
a la tarea de argumentar por qué algunos
filósofos terminan por depravarse (490d); y
la solución exime a la filosofía y culpa a
las condiciones inadecuadas que pudieran
rodear a un hombre potencialmente virtuoso,
volviendo así a la noción bien construida de
\emph{paideía} (491e). Finalmente, una
crítica feroz a los sofistas parece ser el
último trazo de un cuadro donde el filósofo
verdadero se mantiene como tal en un océano
de injusticias y vicios que lo disuaden del
buen camino.

Esta discusión continúa en torno a la
adecuación de la \emph{paideía} en la ciudad
ideal y el papel de la multitud en ella, que
termina por concluir ``que nuestro proyecto,
con tal de ser posible, es el mejor'' (502c).
Este largo camino desemboca en el
descubrimiento de un principio que va más
allá de las virtudes hasta entonces
descritas:

\begin{quote}
``Pero entonces, preguntó, ¿no es lo más
importante lo que hemos visto: la justicia y
lo demás, sino que hay algo mayor aún?''
(504d)
\end{quote}

Presentada como ``la ciencia más alta''
(\emph{mégiston mátheema}) (505a), la idea
del Bien entra en juego ya que ``es por el
recurso a ella como los actos justos y demás
análogos resultan útiles y provechosos''
(505a). Pero esto sólo nos lleva al problema
de definir adecuadamente este concepto.

Rápidamente Sócrates descarta las
posibilidades de que el Bien fuera el placer
o la inteligencia, pero esto es del todo
insatisfactorio para Glaucón, quien termina
por arrastrar a su interlocutor a confesar
que se cree incapaz de comunicar una
respuesta satisfactoria ``con el impulso que
llevamos ahora'' (506e), recurriendo
nuevamente a la analogía.

Esta vez, usando a la vista como ejemplo,
Sócrates nos hace entender a la idea de Bien
como una condición de posibilidad de las
demás virtudes, en particular de la verdad y
del conocimiento (508c). Esto lo vemos
primero en un contexto comunicativo: es el
Bien el que permite el flujo de la verdad
entre el objeto cognoscible y el sujeto
cognoscente (508e). Pero muy pronto,
Sócrates, ayudado de su analogía con el sol,
extiende la noción de Bien a aquello que da
``la existencia y la esencia'' (509b) a los
objetos inteligibles, concluyendo al fin que
el Bien mismo se encuentra ``más allá de la
esencia y la sobrepasa en dignidad y poder''
(509b).

Profundamente exaltado, Glaucón le ruega a
Sócrates que profundice y no omita ningún
aspecto sobre el tema en cuestión, aunque
deba remitirse a la comparación con el sol.
``Muchas cosas seguramente, dije, son las que
he omitido''. Siendo así que una nueva
analogía entra en juego: ahora valiéndose de
un segmento de línea, Sócrates nos hace ver
cuatro secciones que representan a un mismo
tiempo una clasificación ontológica y
epistemológica: sombras-imágenes, cosas
materiales-creencias; cosas inteligibles
hipotéticas-pensamiento discursivo
(\emph{diánoia}), cosas inteligibles no
hipotéticas-conocimiento verdadero
(\emph{nóesis}) (510a-511b). La \emph{nóesis}
es la cumbre de esta clasificación, y ella
misma culmina en el hecho de que al
desprenderse por completo de la hipótesis
para la generación del conocimiento, termina
por ser el único método por el cual se accede
a lo que las cosas son en verdad (ya que la
hipótesis implica la aceptación arbitraria de
algo que no tiene por qué ser cierto) y así,
la vía para llegar ``hasta el principio de
todo'' (511b).

``Te entiendo, dijo, pero no lo bastante
{[}\ldots{}]'' (511c). Pero claro que, aunque
humilde, Glaucón es bastante astuto para
notar y repetir que la posibilidad de
alcanzar la \emph{nóesis} estriba en su
totalidad en el método epistemológico que
permite a la razón tomar contacto con los
objetos por sí misma: la dialéctica.

Libro VII

El libro VII comienza con la conocida
alegoría de la caverna, a través de la cual
Sócrates nos muestra nuevamente la figura del
filósofo, que destaca entre las masas por
tener conocimiento directo de las ideas, y en
particular de la Idea del Bien (representada
por el sol fuera de la caverna) (517b).
Además, Sócrates vuelve sobre la
incomprensión del pueblo frente al filósofo:
``¿no lo matarían si pudieran echarle mano y
darle muerte?'' (517a)

Luego, vemos una fuerte aseveración que da
aún más importancia a la idea del Bien de lo
que teníamos hasta ahora: ``la idea del bien,
que con dificultad percibimos, en el extremo
límite del mundo inteligible, pero que, una
vez entrevista, aparece al razonamiento como
siendo en definitiva la causa universal de
todo cuanto es recto y bello'' (517c), siendo
así que no sólo la verdad y el conocimiento
están íntimamente ligadas a la idea del Bien,
sino que también lo están otras virtudes.

Una cosa que yo quiero colegir es que la idea
de justicia participa de la idea del Bien.
Esto no es explícitamente dicho. Lo más
cercano que tenemos es que, en el mismo
517d-e, Sócrates afirma que no es de
extrañarse que aquel que ha visto la luz
tenga dificultad para deliberar ``en los
tribunales o en otra parte, sobre las sombras
de lo justo o sobre las figurillas cuyo
reflejo son las sombras, y contender sobre la
concepción que de ellos puedan hacerse los
que jamás han visto la justicia misma''. Pero
más bien esto sugiere que la idea de Bien
permite al filósofo ver a la justicia en sí
misma. Yo quisiera ver dicho algo como ``el
Bien da ser a la justicia''. Parece que la
aseveración de que el Bien es causa de todo
lo bello y recto podría resolverlo, pero
tampoco está argumentado que la justicia sea
bella o recta (sobre todo considerando el
argumento de Trasímaco donde la justicia
parecía ser relativa a la fuerza). ¿Podemos
asegurar que el Bien da ser a la justicia?

Un tema que se hace muy presente en este
libro es el de la \emph{paideía}, que nos
devuelve a un plano más tangible. Aquí, este
concepto juega el papel de vía para alcanzar
la idea de Bien en un contexto de aprendizaje
y crecimiento, a diferencia del contexto
epistemológico. Esto devuelve la conversación
al plano social, más propio del diálogo
entero. Se critica nuevamente a los sofistas
(518c) y se reitera la noción de unidad y
cohesión en la polis (519e) antes de dar paso
a la siguiente cuestión a tratar: ``¿No
quieres que a continuación consideremos de
qué manera podrán darse tales hombres, y cómo
se les hará subir a la luz[\ldots{}]?''
(521c) Con lo cual quedan entremezcladas las
discusiones epistemológica y social: ``¿Cuál
podrá ser así, Glaucón, la ciencia
[\emph{mátheema}] que atraiga el alma de lo
que deviene a lo que es?'' (521d)

La primera sospecha es que la \emph{paideía}
hasta entonces construida comprendiera tal
ciencia, pero muy pronto son descartadas
tanto la gimnástica como la música, por
insuficientemente elevadas (522a). Las
siguientes dos propuestas, tentadoras por ser
tan universales, son la del número y la del
cálculo o cuenta (\emph{logismós}) (522c).
Después de una argumentación que apela a la
naturaleza del número como proporción (una
concepción del número muy griega) y como
objeto altamente útil, Sócrates y Glaucón
concluyen que ambas ciencias califican para
la tarea en cuestión; pero no satisfechos con
ello, proponen una tercera: la geometría,
también aceptada por ser ``conocimiento de lo
que siempre es'' (527b).

Interesantemente, las siguientes propuestas
son la geometría de sólidos (58b) y luego la
astronomía (527d). Es interesante porque
desde un punto de vista matemático moderno
(al menos después de Descartes con la
introducción del plano cartesiano) las
geometrías de superficies y de sólidos se
pueden estudiar a través del número. Incluso
la astronomía, entendida como la ciencia ``de
los sólidos en movimiento'' (528e), se ajusta
perfectamente a un espacio cartesiano de
cuatro dimensiones donde la última es
solamente una coordenada (un número) que
representa un instante en el tiempo. Y no
sólo esto, que no es más que una curiosidad,
entrelaza estas cinco primeras propuestas: es
más importante el hecho de que, al ser
ciencias matemáticas, son los ejemplos más
claros de pensamiento dianóico o discursivo.

Dicho esto, no resulta extraño que por más
universal que sea el alcance de estas
disciplinas, no bastan para llegar a conocer
la idea del Bien: ``Todo esto ¿o es que no lo
sabemos? No es sino el preludio de la melodía
que hay que aprender. No creo, en efecto, que
te parezca a ti que los expertos en estas
materias sean dialécticos'' (531d-e).

Así se aborda por primera vez el concepto de
dialéctica en el este libro, y a continuación
no se hace una explicación exhaustiva ni de
qué es en verdad la dialéctica ni de por qué
es decisivamente la ciencia de la esencia.
Sócrates tan sólo aclara que paso a paso
(según yo haciendo alusión a la noción de
\emph{hypothésis} entendida como peldaño y no
como axioma, cosa que explicó en el libro
anterior cuando caracterizaba la
\emph{nóesis}) y renunciando del todo a los
sentidos se alcanzará el bien en sí mismo,
siendo ésta la ``melodía de la dialéctica''
(532a).

Sin más, el libro da un último giro hacia el
carácter práctico de la diálectica: ``Lo que
queda ahora, proseguí es regular la
distribución de estas enseñanzas: a quiénes y
de qué manera.'' (535a). Vemos una vez más
las características virtuosas de aquel que
estudiará la dialéctica (535d, 536a) y se
reitera la importancia de las ciencias como
la aritmética y la geometría, ahora ``como la
propedéutica de la dialéctica'' (536d). El
libro termina haciendo un breve comentario
acerca las posibilidades de existencia que
tiene la ciudad discutida: ``Si su
realización es difícil, no deja por ello de
ser posible'' (540d).

Creo que hay al menos dos temas en La
República que sería muy lindo discutir: si la
polis ideal presentada en el texto podría o
no existir, y si el método discursivo que usa
Platón es \emph{nóesis} o \emph{diánoia}. Y
bueno, creo que son parte de un mismo
problema. Si la polis no puede existir (que
en mi opinión es por mucho lo más probable),
entonces estamos obligados a decir que el
discurso es del todo dianóico (por estar
basado en un supuesto no existente), lo que
resulta extraño dado que, por un lado, Platón
escoge el diálogo (que al menos
etimológicamente es casi como decir
\emph{diálektos}-discusión) para explicar
todo esto, cosa que sugiere fuertemente el
uso de la dialéctica para generar el
conocimiento; y por otro, porque establece
explícitamente que la diánoia no es el método
predilecto para el conocimiento.

En última instancia, todo se reduce a
preguntarnos cuáles son las pretensiones de
Platón en este diálogo: ¿qué quiere hacer ver
a sus lectores? Definitivamente no a la Idea
del Bien, ni tampoco por qué es la dialéctica
el método por excelencia para la nóesis. En
realidad, nos debe quedar claro que éste no
es un tratado sobre ontología ni
epistemología, y en ese sentido Platón no
está obligado a profundizar sobre estos temas
ni fundamentarlos del todo; cosa que se hace
evidente al menos cuando Sócrates explica a
sus interlocutores que no está posición de
transmitirles la idea del Bien con la
construcción que han logrado hasta
entonces--- ni hasta el final del diálogo.

Pues bien, entonces, ¿qué sí pretende hacer
ver a sus lectores? Y, ¿por qué es necesario
abordar temas tan abstractos para
conseguirlo? La República, Politeia, es al
fin una discusión en torno de la idea de
justicia, que vemos manifiesta como principio
tanto humano, en el alma, como social, en la
polis. ¿Será acaso que Platón no se puede
contener para hablar sobre su teoría de las
ideas? Más bien parece que la idea del Bien
surge naturalmente, y acaso debería surgir
naturalmente en toda investigación
suficientemente exhaustiva respecto de la
justicia, puesto que ésta se encuentra por
completo sujeta a aquélla, que casi ha
llegado a ser Dios.

Resulta ser que Sócrates, en el libro II,
emprende al menos dos tareas: primero,
defender la justicia por sí misma y no por lo
que causa (358b), y luego, responder a las
aseveraciones de Trasímaco en las que
argumentaba que es siempre mejor la
injusticia que la justicia (343e-344c), y que
la justicia es lo beneficioso para el más
fuerte (338c).

Pues bien, el papel de la idea del Bien es
fundamental para ambas tareas: el situar a la
justicia en el contexto ontológico al que se
ha llegado en el libro VI le confiere
perfecta justificación como ente en sí mismo:
es una idea, y por ello, no sólo no es
necesario observar las consecuencias que se
desprenden de la justicia, sino que éstas son
del todo prescindibles al ser imperfectos
entes del mundo material.

Y como respuesta a Trasímaco, toda
argumentación de las que vemos en el primer
diálogo pierde por completo sentido si ya
hemos llegado a ver que la justicia no es
sino una manifestación, o un ente que
participa directamente de la idea del Bien.
Es decir, después de trazar esta intrincada
relación entre la idea del Bien y la
justicia, resulta imposible argumentar que lo
injusto fuera mejor que lo justo, al ser ya
lo justo cosa tan buena. Y que la justicia
fuera el beneficio del más fuerte queda
también descartado porque aquello supone la
desgracia del débil, cosa antinatural en un
idea buena en sí misma.

Parece ser que todo termina perfectamente
bien justificado con la idea del Bien. Pero,
¿cómo justificar la existencia de la idea del
Bien? Esto con toda intención de demostrar
que esta idea no es tomada como hipótesis o
axioma en la teoría platónica, cosa que
derrumbaría por completo la parte ontológica
y epistemológica. Además, esta es una
pregunta importante dado que, según entiendo,
éste es el único diálogo en el que Platón
aborda la idea del Bien.

La única respuesta que yo encuentro a esta
pregunta es que toda la filosofía platónica,
al girar en torno de las virtudes y las
ideas, gira en torno de la idea del Bien y
así, conforma, en su totalidad, esta
demostración. Quiero decir que las virtudes
son virtuosas en tanto que son buenas, y en
ese sentido hablar del Bien es sólo dar un
nombre a una característica común y esencial
de las virtudes. Decir que el Bien exista es
tanto como decir que la filosofía platónica,
entendida como la filosofía dirigida al
cultivo de lo virtuoso, exista.

¿Es muy descabellado? A mí me da mucho la
impresión de que la filosofía griega en
general gira en torno de la virtud, sea
conocimiento, belleza, templanza, etc.
Seguramente es arbitrario y no del todo
correcto decir que la filosofía griega o la
platónica es la filosofía dirigida al cultivo
de lo virtuoso, pero al menos esta
interpretación es muy congruente con el hecho
de que lo más cercano a un dios sea el Bien
(y que el Dios platónico sea tan bueno). En
fin, estas refelxiones se desprenden casi por
completo de la República y, si sí es
descabellado todo esto, quizás haya algo de
bibliografía para ver otra cara de la
filosofía griega.

\section{La República: de la democracia a la
tiranía}
\label{section-republica-democracia}

The final essay of the course, of December
2017: ``Lectura y análisis sobre la
corrupción de la democracia en La
República''. It concerns Book VIII, which
the essay calls Book XIII, and the passage
it cites as 163a--164a is 563a--564a.

En el libro XIII de La República, Sócrates se
propone hacer una revista de los cuatro tipos
de gobierno que, según él, hacen falta para
tener clasificadas todas maneras de regir una
ciudad. Esto lo hace tomando muy en cuenta la
analogía polis-anthropos, que es el supuesto
que le permite investigar el alma humana al
mismo tiempo que la estructura social en una
ciudad. Habiendo ya examinado con gran
detenimiento el caso de la aristocracia, y
clasificádolo como bueno y justo (544e), los
restantes son la timocracia, la oligarquía,
la democracia y la tiranía.

En su investigación, Sócrates describe cómo
es la transición de un Estado a otro,
comenzando por la aristocracia y prosiguiendo
en el orden mencionado. Entender cada
transición supone la comprensión de las
debilidades fundamentales de cada
Estado-hombre, de tal forma que se exhibirán
las virtudes y los vicios de cada caso. En
este ensayo trataré de explicar la transición
de la democracia a la tiranía haciendo un
recorrido por la construcción de la primera y
analizando la forma en que Platón argumenta
la superioridad de la aristocracia sobre
ella.

Para entender en qué consiste la corrupción
de la democracia, es necesario analizar
primero cómo fue su surgimiento y en qué
consisten sus principales virtudes. El Estado
que le precede, la oligarquía, nos es muy
familiar porque en muchos sentidos lo
observamos en la modernidad: una ciudad
regida por una minoría que no procura el
bienestar del pueblo sino la riqueza propia.
Sócrates explica cómo de tan tremenda
desigualdad surge una sublevación del pueblo,
y de ella, la democracia: ``Nace, pues la
democracia, creo yo, cuando los pobres,
victoriosos de sus contrarios, matan a unos,
destierran a otros, y comparten
igualitariamente con los que quedan el
gobierno y las magistraturas, que en este
régimen, además suelen cubrirse por sorteo''
(557a).

Así, las dos virtudes principales de la
democracia serán la libertad y la igualdad.
En un inicio, ``puede parecer por extremo
bella esta república [\ldots{}]'' (557c),
pero Sócrates no se deja llevar por una falsa
impresión y rápidamente comienza a delinear
el inminente caos que se desprende de tener
por libertad al principal derecho: ``no hay
ninguna obligación de gobernar, ni aun para
quien sea capaz de hacerlo, ni a la inversa
de obedecer si no te agrada, ni de hacer la
guerra cuando la hacen los otros, ni de
observar la paz cuando los demás la observan,
si no quieres paz'' (557e).

Platón entiende la democracia como el paso
natural de un sistema cuyos líderes no se
preocupan más que de sí mismos a otro cuyos
miembros repudian la noción de autoridad; y
por ello, para mostrar que la democracia es
insostenible, es necesario explicar por qué
la falta de un orden jerárquico lleva
inminentemente a la corrupción de la sociedad
en la polis y del alma en el hombre. Además,
ésta es una preocupación mayúscula para la
filosofía platónica, pues, al menos en La
República, es evidente que nuestro filósofo
ordena los objetos que estudia y siempre
distingue entre ellos uno supremo: el logos
en el alma, el rey filósofo en la polis, las
ideas entre los entes, el bien entre las
ideas. Muy pronto escuchamos: ``¡Con qué
arrogancia se pisotean estas máximas
{[}\ldots{}]!'' (558b).

En lo que respecta al alma democrática,
Sócrates nos hace ver cómo el joven
demócrata, lleno de odio hacia su padre
oligarca, destierra de su alma al pudor, a la
templanza y a la moderación, y en su lugar
satisface ``una multitud de deseos
superfluos'' (560d): es la epithymía la parte
del alma que se corrompe. Esto resulta en una
confusión absoluta de las demás virtudes, y
de los mismos placeres, negándose el
demócrata a aceptar el valor de ciertos
placeres sobre otros (561c). ``No hay en su
vida, en suma, ni orden ni disciplina''
(561d).

En este sentido, podemos ver en el alma
democrática que el desorden que la
caracteriza está en estrecha relación con la
supremacía de una parte del alma que
realmente debería ser subordinada de otra (la
epithymía del logos). Es decir, Platón no
sólo quiere argumentar la necesidad de un
principio que rija a los demás, sino que nos
debe decir cuál es; o más bien, en este caso,
cuál no es.

Vemos ahora cómo lo hace. Creo que hay al
menos dos puntos principales en la
argumentación. El primero se puede resumir en
una cita que apareció desde que comenzaba la
descripción de la democracia: ésta
``[\ldots{}] dispensa indistamente una
especie de igualdad tanto a los iguales como
a los desiguales.'' (558c). Sócrates
desarrolla esta idea en 163a-164a partiendo
de un hijo que se iguala a un padre hasta
llegar a un caos absoluto donde parece no
haber distinción entre cualesquiera
individuos e incluso entre hombres y
animales. Esto culmina con el hecho de que
``todo exceso en el obrar produce
generalmente, como compensación, un cambio
considerable en sentido contrario'' (563e), y
de ahí que de la máxima libertad se origine
la máxima esclavitud, la tiranía.

El otro punto que compone la argumentación
apela a una desidealización de la democracia:
así como la caída de la aristocracía suponía
una corrupción natural en el hombre y en la
polis, debemos también admitir que no es
posible una ciudad que carezca absolutamente
de dirigentes. Sócrates hace una analogía con
un panal de abejas, donde los dirigentes del
pueblo, aunque no son alados como los propios
de la oligarquía, tienen aguijón, mientras
que los dirigidos, por así decirlo, no.

Más aún, nos es presentada una división
tripartita de las clases sociales: los
dirigentes, los ricos, que han obtenido su
ganancia por ser más ordenados por naturaleza
(564e), y los restantes, la mayoría que no
posee mayor riqueza (565a). Por cierto que
esto no tergiversa la esencia de la
democracia, al menos en tanto que esta última
clase ``es el grupo más numeroso y el que,
cuando se reúne en asamblea, es el soberano
en la democracia'' (565a).

En este escenario, sucede lo siguiente (565b-
): la clase dirigente despoja a la clase rica
de su fortuna para repartirla al pueblo en
virtud de la igualdad que define a la
democracia, lo que genera cierto rencor en
esta clase de potenciales oligarcas, u
oligarcas declarados, siendo así que se
desata una querella. Entonces, asustado, el
pueblo busca alguna ayuda que le proteja ante
semejante sublevación de los simpatizantes
del antiguo régimen, y surge así un falso
amigo que se convertirá en tirano.

El diálogo continúa haciendo la descripción
de cómo el futuro tirano goza de la simpatía
del pueblo pero no es capaz de controlar su
hambre de poder, hasta convertir su propia
existencia en mera esclavitud y sufrimiento
disfrazado de placer.

Hemos visto, pues, cómo Sócrates argumenta de
dos formas distintas el colapso de la
democracia. Creo que el primero es el que
apela de manera más directa a lo que define
la democracia, en el sentido de que la raíz
del caos es la misma igualdad. El recurso
fundamental que utiliza Platón para
evidenciar la insostenibilidad de la
democracia es la exageración del principio de
igualdad. ¿Qué tan válido es hacer esto? En
mi opinión, es admisible porque lo que aquí
intenta exhibir Platón es, esencialmente, que
la igualdad no se sostiene sola, o al menos,
que no es capaz de sostener a los demás
principios o virtudes como subordinados de
ella. Esto contrasta enormemente con la
justicia o incluso la idea del bien, en el
siguiente sentido: ¿qué pasaría si exageramos
de manera equivalente una ciudad cuya máxima
fuera la justicia o el bien? Parece ser que
la argumentación de los libros anteriores
demuestra que al contrario de llevar el
Estado a su disolución, aquella situación
resultaría en una polis ideal. Por esta
razón, entiendo que la exageración no va más
allá de un recurso discursivo que me parece
lícito en tanto que el objeto de estudio es
la igualdad como principio supremo y no la
posibilidad de una ciudad absolutamente
igualitaria.

La segunda manera de argumentar empata
perfectamente con esta idea. El delicadísimo
juego entre lo hipotético y lo posible, que
permea a lo largo de toda la República, se
hace evidente una vez más cuando Sócrates
reconoce que el Estado democrático está
dividido en clases sociales. Creo que esta
parte del argumento está situada en segundo
lugar porque su estructura está más
relacionada con el modelo de tiranía al que
pretende desembocar. En un plano más
práctico, Platón explica cómo una discordia
tan natural como la que le fue necesario
reconocer como causa de la crisis del Estado
ideal convertirá la democracia en un régimen
que vuelve a ser autoritario como la
oligarquía, pero nunca podría volver ser ella
misma porque el fundamento de la democracia
es precisamente la negación de aquella
estructura, y por lo tanto, termina en un
autoritarismo distinto que resulta ser mucho
más radical y violento, y que llama tiranía.

Por fin, Platón parece haber argumentado
ampliamente la necesidad de un principio
distinto del de la libertad o la igualdad,
que se declare superior a éstos y que sea la
máxima que rija la polis. Así, un repudio
desmesurado a la autoridad es insostenible en
una sociedad: es necesario que el pueblo
reconozca alguna entidad superior a él que
tenga licencia de privarle de ciertas
libertades para asegurar la prosperidad de la
polis (que esta entidad sea o no el rey
filósofo, no está argumentado en esta parte
del diálogo). Esto reitera la necesidad de
una jerarquía social, que se corresponde con
una jerarquía en el alma. Ahora bien, el caso
del hombre democrático nos muestra cómo
sucumbir ante los placeres, necesarios o no,
no es la vía que llevará a una buena vida,
descartando así a la epithymía como la parte
del alma que debería regir sobre las demás.

\section{Notas de lectura sobre Aristóteles}
\label{section-notas-aristoteles}

Reading notes of February and March 2018,
written while reading for Filosofía del
Lenguaje (\emph{De interpretatione}, the
\emph{Metaphysics}, the \emph{Physics}) and
for Historia de la Filosofía II (the
\emph{Nicomachean Ethics}, Epicurus). They
are rough notes, as written; page numbers
refer to the editions I was reading.

\medskip\noindent
\emph{Sobre la interpretación.}

``La confusión platónica del ser con un
predicado general'' y luego ``es quizá esto
lo que mejor diferencia la ontología
aristótelica de la platónica''. ¿Qué quiere
decir esto? Que Platón entendía en el verbo
ser algo más que un artefacto gramatical que
nos permite ``decir la afirmación o negación
de algo''. ¿Qué entendía Platón en el ser?

¿Qué hace el verbo para Aristóteles?
Cosignifica- prossemaínon

Aquí no se va a discutir por qué hay
biyección entre nombres y cosas.

La argumentación al principio de no
contradicción: no (no-contradicción) implica
``cabría afirmar todo lo que uno negara y
negar todo lo que uno afirmara'', ie. no
existe la verdad.

\medskip\noindent\emph{Metafísica.}

Acaso la mayor virtud de este capítulo es
hacer al lector reconsiderar la utilidad de
la filosofía\ldots{} nula.

Básicamente le preocupa argumentar que la
ciencia que estudia al ente en cuanto ente es
sólo una.

``y es la misma ciencia la que contempla no
sólo la sustancia, sino los atributos'' pg163
Eso fue el cap 2.

Cap3. Todos los entes se sirven de axiomas

El principio supremo no debe ser hipotético
(non conditionale) y el más conocido, razón
por la cual el principio supremo es\ldots{}

LA TRAGEDIA MÁXIMA:

Todos piensan que el más alto saber es el
saber de los principios.

El principio supremo es el principio de no
contradicción. pg167

Y así fue como la Edad Media quedó condenada.

\emph{hadynanton ama ypolambenein ton ayton
einai kai me einai to ayto}

imposible admitir que una cosa es y no es.

``el axioma de axiomas''

La respuesta aquí es: la lógica no es el
único modo de adquirir conocimiento. Sobre
todo en el sentido en el que la filosofía es
conocer para vivir, y conocer por conocer,
hay que comprender los límites del logos y no
enfrascarse en ``demostraciones infinitas''.

Dice él ``de algunas cosas no se debe dar
demostración'' pg 169. En efecto, pero esto
implica que necesitemos un axioma para
justificarlas todas.

Es un problema al fin ligüístico, de
lenguaje. Sí porque como las cosas están
todas mal definidas, ninguna acerción puede
ser terminantemente verdadera o falsa. La
misma noción de verdad, que al fin subyace al
principio de no contradicción, está mal
definida. Empeñarse como lo hace Aristóteles,
en la persecución por la ciencia máxima, en
dar fundamento lógico a la filosofía, es un
error. La filosofía no persigue fundamentos
lógicos, Sócrates lo sabe.

Veo dos opciones en la vida. Divagar libre y
felizmente, asumiendo que uno no sabe
realmente, o hacer matemáticas suponiendo
explícitamente los axiomas (como a
Aristóteles ni a nadie le gusta por son
axiomas hipotéticos) y divertirse mucho
haciendo eso de juego. En fin, que la vida es
un juego, es claro. Y yo aquí tratando de
determinarlo con certeza.

La discusión: pg 171, ``es evidente que esto
al menos es verdadero en sí: que la expresión
``ser'' o ``no ser'' significa algo
determinado''. Cómo se atreve a decir eso. Mi
opinión no implica que nada exista\ldots{}

162 ``a cada concepto se le podría imponer un
nombre diferente''

174 La esencia de la cosa es la esencia del
nombre.

176 Qué bien lo dice ``todos los que dicen
esto destruyen la esencia y la sustancia''.

El argumento letal 187: ¿por qué no te tiras
al precipicio? Pero eso es reducción al
absurdo! fútil sin el principio de no
contradicción. Cómo un razonamiento me va a
tirar al precipicio.

\medskip\noindent\emph{Física.}

\emph{diorisménon}

\emph{synechés}

Las cosas existen por naturaleza tienen
tendencia al cambio y las cosas que existen
en virtud de otras causas no.

``La naturaleza es un principio y una causa
del cambio y del reposo de aquella cosa en la
que se da primariamente por sí misma'' o bien
``es la primera materia que subyace a cada
cosa de las que tienen en sí mismas el
principio del movimiento y del cambio''.

``La naturaleza es la figura y forma según la
deinifición'' Dada una definición, la
naturaleza es h' morfos kai to eidos to kata
ton logon.

``La naturaleza concebida como proceso de
generación es un camino hacia la
naturaleza''. El fin de la naturaleza es ella
misma

pg 29 dos artes dominan la materia (hyle) la
que usa y la que produce (texne productora).
La que usa es conocedora de la forma y la que
produce lo es de la materia.

Cap3

``Se llama causa a aquello a partir de lo
cual se genera algo existente''

``Se le llama causa a la definición (logos)
de la eidós.

``Se le llama causa. al comienzo del cambio''

``Se le llama causa al telos, propósito''

``Se le llama causa a todo aquello que se
genera entre la cosa y el fin'' (dani: Todo
lo previo al fin en tanto que causa el fin)

Una causa de algo puede ser una propiedad
accidental de otra causa. Causa es amplísimo.

pg33 (23) busca la causa más exacta de cada
cosa (dani: no existe)

Una causa accidental existe por suerte o
azar.

Está diciendo: algo pasa por una causa bien
determinada, y en consecuencia de ese suceso
se da otro. El segundo ocurre ``por suerte''
y acaso la causa del primero lo es también
del segundo.

En conclusión: la suerte no es causa última
de nada (dani: en ese sentido no existe? yo
quiero que aristóteles no quiera que exista
la suerte)

Que hable de la infinitud de la causas: vaya
no numerabilidad! vaya complejidad.

La suerte ligada con lo bueno y lo malo; la
suerte ligada a la acción (praxis) de los
seres animados. La diferencia está en la
posibilidad de elegir.

``el caballo vino por azar, porque se salvo
al venir, pero sin venir para ser salvado''

Azar = automaton: cuando algo se da él mismo
en vano. Por sí mismo vs por naturaleza

El Nous y la fisis anteriores al azar y a la
suerte.

Las cuatro causas: la materia Hyle, la forma
Eidós, lo que inicia el proceso Kinesan, y el
propósito `Eneka. A ellos se pregunta ``por
qué?''

\medskip\noindent\emph{Ética a Nicómaco.}

Prólogo en defensa de Aris: el logos de la
pragmateía (escribir como escribe aris) es
abierto y libre y dialéctico porque su
publicación era la lectura, ékdosis.

Dani: ¿será que de la existencia del lenguaje
y el lógos se desprende la necesidad de la
ética? de fundamentar la ética.

La sofística se fundamenta por un lado en el
relativismo

Mi queja fundamental a Aris: la esencia de
las cosas no está bien definida; la
abstracción, la generalización, son siempre
injustas. Y otra: lo de que la filosofía es
la búsqueda de causas esenciales, de las
causas más generales. Y otra: si la ética es
búsqueda de felicidad (y la filosofía ética
la labor de definir esto bien), él dice que
la más feliz es la vida contemplativa.

De Platón a Aristóteles sobre el Bien en sí.
Las cosas buenas no lo son en tanto que
participan de la Idea del bien, lo son en
tanto que se ajustan al fin (dado
naturalmente, esencialmente) de cada
actividad. La Idea del bien se desprende de
la analogía de los casos particulares y en
este sentido el ``Bien trascendente'' es
substancia y no idea.

No es la preservación de la vida la finalidad
de la vida humana porque la preservación de
la vida ya es la finalidad de la vida de las
plantas y aquí buscamos lo propio: la
finalidad del hombre en tanto hombre. Esto es
una barbaridad de razonamiento porque el
hombre en tanto hombre va a definirse de
acuerdo a un parámetro arbitrario, la razón.
Además, ¿unicidad de la finalidad de la vida
humana?

Lo que tiene de bueno Aris es que dice que
las abstracciones, las generalidades, salen
del mundo, salen de las cosas, de las
particularidades. Por acto del intelecto
humano. ¿Cuánta realidad hay en el intelecto
humano?

¿Por qué el bien del hombre tiene que ser un
fin en sí mismo?

Parece ser que Aristóteles define el término
medio sólo para ciertos conceptos. Hay cosas
que pueden hacer mucho o poco, en exceso o en
defecto, y cosas que no son susceptibles al
exceso o al defecto, por ejemplo, el mismo
término medio.

Respecto a todo esto: nunca estamos realmente
en un extremo, ni realmente en medio. Si nos
imaginamos una línea, yo me imagino que nunca
estamos en un punto dado. Que al ser éstas
cualidades del alma, de la acción, que el
mismo Aristóteles reconoce inaccesibles a la
razón, la analogía de la línea es peligrosa,
y sólo nos sirve para aproximar nuestras
cualidades. Hay que entender que son
dinámicas, que no están bien definidas, y que
la filosofía, la reflexión sobre ellas, es
más una guía abstracta que debe desembocar en
acciones.

\medskip\noindent\emph{Epicuro.}

Siempre que se habla sobre el ser y el no
ser, parece haber una especie de reducción al
absurdo. Pues es que el mismo concepto de no
ser es producto de una operación lógica (como
la reducción al absurdo). No es que al
reducir al absurdo estamos aceptando la
lógica, al plantear el problema por el no ser
ya se ha aceptado este esquema.

\section{El término medio en Aristóteles}
\label{section-termino-medio}

An answer of March 2018 for Historia de la
Filosofía II to the question ``¿En qué
consiste el término medio para
Aristóteles?'', on the
\emph{Nicomachean Ethics}.

Una primera definición del término medio es
geométrica: ``Llamo término medio de una cosa
al que dista lo mismo de ambos extremos''.
Pero Aristóteles rápidamente se separa de un
modelo demasiado matemático: ``pero el
término medio relativo a nosotros no ha de
tomarse de la misma manera'': como el tema es
la ética, el término medio lo será respecto
del exceso y el defecto, siendo éstos los
extremos de lo que se nos figura un continuo
o un orden de acciones o actitudes respecto a
cierta cosa.

Aristóteles nos da bastantes ejemplos de
estos continuos en los que hay extremos y
término medio: honor, miedo, ira, verdad.
Para cada uno, el autor nos dice cómo se
llama al hombre que se caracteriza por
hallarse en uno de los extremos o en el
término medio. Siempre, la finalidad de esta
clasificación es hacer ver que el término
medio coincide con lo virtuoso y lo bueno.

Algo interesante es que, para Aristóteles, no
de cualquier cosa se puede plantear este
modelo de extremos y término medio, por
ejemplo, del mismo término medio: no hay
exceso ni defecto de término medio. Sobre las
mismas líneas, una excepción a la regla de
que lo bueno es el término medio, es
precisamente la virtud, puesto que ``con
respecto a lo mejor y al bien, es un
extremo''.

Como sea, creo que no vale la pena fijarnos
en estas excepciones que podrían llevar a
alguna aporía en el razonamiento de
Aristóteles, y más bien atender al mensaje
ético que trata de comunicarnos. El término
medio es un recurso que aunque es fácil de
entender mediante un modelo geométrico, se
refiere actitudes humanas, y conduce a la
moderación. Aristóteles trata de comunicarnos
que el bien y lo virtuoso se encuentran entre
los extremos nuestras actividades: debemos
aprender a reconocer cuáles de nuestras
acciones son susceptibles a pervertirse hacia
uno u otro extremo, entender en qué consiste
el punto medio entre ellos y esforzarnos para
alcanzarlo. Además, está en nuestras manos
movernos hacia un lado u otro del continuo, y
seguramente encontrar el término medio será
lo más difícil y cansado. Veo en el discurso
mucho énfasis en los conceptos de autognosis
y autocracia: ``Debemos, por otra parte,
tomar en consideración aquellas cosas hacia
las que somos más inclinados (pues unos lo
somos por naturaleza hacia unas y otros hacia
otras)''.

En conclusión, creo que el término medio para
Aristóteles es un recurso que se ilustra
claramente como el punto medio en una recta
con dos extremos que representan el exceso y
el defecto, pero no tiene otra intención que
la de promover la moderación como medio para
lograr la virtud y el bien, en un continuo
que realmente no está definido matemática
sino prácticamente (``todo lo que se diga de
las acciones debe decirse en esquema y no con
precisión'') en el contexto de la vida diaria
del hombre.

\section{El sabio cínico según Luciano}
\label{section-sabio-cinico}

An answer of May 2018 for Historia de la
Filosofía II, on Lucian's dialogue
\emph{El cínico}.

El sabio cínico es una figura que, en primer
lugar, llama la atención por su aspecto:
``estar sucio, ser peludo, tener capote,
llevar el pelo largo, ir descalzo''. Se trata
de una persona que vive en condiciones
mínimas para la supervivencia. Ahora bien, le
llamamos sabio a cuenta de que su decisión de
vivir de esta manera estriba no en la locura
sino en la elevada capacidad de
discernimiento que le permite notar que el
bienestar no radica en la satisfacción del
placer sino en el control de la avidez.

Vemos en el diálogo El cínico un discurso de
crítica fuerte a esta avidez, ambición,
hedonismo, que pudo encontrar Luciano en los
hombres de su época tanto como podemos
hacerlo nosotros con los de la nuestra y, hay
que decirlo, en nosotros mismos. La figura
del cínico es el caso límite de una persona
que se ha librado de estos vicios para
abrazar de lleno el camino contrario. Vemos
en él una persona que actúa de acuerdo a sus
palabras con suficiente valor como para
resistir el ser comparado con un perro o un
mendigo.

Mi reflexión respecto a los cínicos es la
siguiente. No vemos en esta escuela una
propuesta filosófica realmente novedosa.
Habíamos visto los principios que rigen su
estilo de vida, por ejemplo, en Sócrates,
conocido por ser pobre en el campo práctico,
e ignorante en el campo teórico (su gran
sabiduría radicaba en la consciencia de su
ignorancia). Los cínicos son una
manifestación más de la idea ampliamente
abordada en la filosofía griega de no
intentar a toda costa alcanzar un objetivo
que en primera instancia se nos figura
deseable, sino cuestionar el verdadero valor
de dicho objetivo para notar que resulta más
benéfico (aunque quizás sea más difícil)
dejarlo a un lado para encontrar el bienestar
en lo que ya se tiene. En efecto, es una
tarea nada trivial la de resistir una vida de
mendigo para el cínico, o la de aceptar la
muerte para Sócrates. En cualquier caso, la
reflexión filosófica, que en última instancia
estriba en la capacidad de discernimiento y
en el logos, es la herramienta esencial
utilizada por Sócrates, Platón, los cínicos y
los estóicos, entre tantos otros, para llevar
a cabo la interminable lucha contra todas
aquellas cosas que nos alejan del bienestar o
eudaimonia.

Sin embargo, no se trata aquí de echar a un
lado la corriente de los cínicos. El hecho de
que los principios esenciales de su enseñanza
no nos parezcan realmente novedosos no es
razón para olvidarlos. La existencia de esta
corriente, junto con la estóica, resulta ser
un factor histórico clave que nos hace notar
el impacto de las ideas socráticas y
presocráticas a través del tiempo, y, en fin,
era la sustancia que constituía la filosofía
en de aquel momento.

\section{Language as an object of reflection}
\label{section-language}

These are the general questions about
language as I met them in the course
Filosofía del Lenguaje (UNAM, Facultad de
Filosofía y Letras, semester 2018-2, taught
by Erika Rebeca Lindig Cisneros). The course
approached the philosophy of language
through the history of its problems, with
emphasis on the exchange between philosophy
and rhetoric. Its programme recalls that
reflection on language became central to
philosophy only in the twentieth century,
after what Rorty called the linguistic turn
and Foucault the return of language. What
follows is the outline of the course with
what I wrote down while reading. The remarks
marked as mine are a student's notes of
2018, not established claims. My essays for
the course are in
Sections~\ref{section-farmacia-de-platon}
and \ref{section-derrida-juego}.

\medskip\noindent
\emph{Rhetoric and philosophy.}
The first unit read Plato's \emph{Cratylus}
and \emph{Phaedrus}, Derrida's ``La farmacia
de Platón'', Aristotle (\emph{De
interpretatione} 16a--17a, \emph{Rhetoric}
1354a--1359b, \emph{Topics} 100a--101b) and
Heidegger's ``El camino al habla''.

In the \emph{Phaedrus} writing is offered as
a \emph{pharmakon} for memory, a word that
means both remedy and poison. Derrida's
reading, as I understood it, is that the
\emph{pharmakon} opens the very possibility
of the oppositions true and false, good and
bad, remedy and poison, without finding a
place inside them; and that the weakness
blamed on writing is already inscribed in
living memory, which is finite.

Aristotle's scheme at the beginning of
\emph{De interpretatione} (16a3--8) relates
four things: written marks are symbols of
spoken sounds, spoken sounds are symbols of
affections of the soul, and these affections
are likenesses of things. Writing and speech
differ among peoples, while the affections
and the things are the same for all. A verb
moreover \emph{co-signifies} time (16b6). My
notes add that Aristotle does not discuss
here why names and things should correspond
one to one.

\medskip\noindent
\emph{Language and reason in the moderns.}
The second unit read Locke's \emph{Essay
Concerning Human Understanding} with de
Man's critique of it, and Rousseau's
\emph{Discourse on the Origin of Inequality}
and \emph{Essay on the Origin of Languages}
with Derrida's ``La lingüística de
Rousseau''.

The programme calls Locke's position
instrumentalism: words are instruments for
ideas that come from experience. From his
argument against innate truths I kept this:
truth is acquired together with signs, which
cannot be separated from the possibility of
being true or false. A general term captures
less the things it covers than the property
by which they are associated.

Rousseau writes that general ideas can enter
the mind only with the help of words.
Derrida looks for the circles in his account
of the origin of language. If men need
speech in order to learn to think, they need
even more to know how to think in order to
find the art of speech; and one cannot say
whether society, already united, instituted
languages, or languages, already invented,
established society. For Derrida this
logical circle also marks out a field of
objects of its own, that of a science of
language.

\medskip\noindent
\emph{The linguistic turn.}
The third unit went from German Romanticism
(Hamann, Herder, Humboldt) to Nietzsche's
\emph{Sobre verdad y mentira}, with pages of
Gadamer's \emph{Verdad y método}.

From Humboldt I noted that language is born
of itself, and that the differences among
languages come from the differences in
spirit among peoples. From Lafont's reading
of Heidegger: language does not merely
express in words what is already unveiled,
but is responsible for a being's appearing
as such; what things are becomes what is
open to us in a given historical language.

\section{Sobre La farmacia de Platón}
\label{section-farmacia-de-platon}

A text of February 2018 for Filosofía del
Lenguaje, titled ``Ensayo sobre el
lenguaje'', written after reading the
fourth chapter of Derrida's ``La farmacia
de Platón''.

Esto es lo que escribí cuando terminé el
capítulo y creo que no captura nada mal mi
idea:

Conforme avanza el capítulo cuatro, Derrida
nos desliza de la escritura al habla, de la
hipomnesis a la mnesis, de la magia al logos.
Comenzamos por ver cómo la escritura se
encuentra por debajo de la memoria como el
temible fármacon que nos alejaría por siempre
de la memoria, de la verdad. Pero la manera
de decirlo, a través del fármacon, condena el
discurso. ¿Qué quiere decir que la escritura
sea un fármacon de la memoria? ¿Qué gravísima
implicación vemos en este hecho? Derrida nos
hace ver que el mero planteamiento de este
problema, la dicotomía escritura-memoria,
falsedad-verdad, en cuyo corazón se encuentra
el fármacon, definido a un tiempo como
remedio y veneno, exhibe una fractura
fundamental en la memoria que no sólo exime a
la escritura del daño que podría hacerle,
sino que la iguala; a través del mito
egipcio, y de todas las imágenes que se
sobreponen o acaso se hilan, terminamos por
ver en Sócrates un brujo y en el logos un
hechizo.

El logos está roto. Siempre lo ha estado, e
ignorar eso es vivir engañado, tan engañado,
sí, como si no lo ignoraras (pues no
ignorarlo es aceptar la imposibilidad de la
sobriedad) pero sin reconocerlo. Y en ello
está el piso de este asunto. No es la última
palabra ---no existe la última palabra, el
discurso se pliega sobre sí mismo
infinitamente y por eso hemos abandonado la
empresa de desenredarlo; y por lo mismo el
discurso existe, nos encontramos
inevitablemente sosteniéndolo: sostenemos
explícitamente la paradoja, reconocemos no
saber nada salvo eso mismo; vemos en Sócrates
un brujo y en el logos magia.

El martes yo trataba de explicar estas ideas
a través del concepto de la memoria porque,
no sólo en la lectura de Derrida sino que
también en la de Detienne, se hace muy clara
la asociación entre memoria y verdad, de tal
forma que el fármacon, cuando es presentado
como cura para la falta de memoria, adquiere
también el significado de cura para la falta
de verdad. Así que, en fin, el problema tiene
que ver con la verdad.

Creo que la propuesta de Derrida es que la
verdad se encuentra fracturada desde su
definición, misma que implica la posibilidad
de falsedad. Que exista la falsedad quiere
decir que existe una frontera del
conocimiento, y que hay quienes se equivocan
respecto de ciertas cosas. Surge una
pregunta: ¿cómo saber que se sabe? Es curioso
que Sócrates sea llamado el más sabio de los
hombres por ser el único en reconocer que no
sabe; anécdota que va de la mano con lo que
Derrida llama ``mandar de paseo a los
mitos'', o ``jairein'', suceso en el que
Sócrates no se queda con una ni otra
explicación del mito de Oritia.

¿Cómo saber cuál de las explicaciones es la
verdadera? Pero, mucho más: ¿qué efecto tiene
sobre nosotros la existencia de esta pregunta
tan difícil de responder? La última pregunta
que hice en la clase, y que hice mal
porque dije Detienne en vez de Derrida, era
por la historia de este problema. Siempre han
existido explicaciones, digamos,
irracionales, de las cosas. Siempre ha
existido la tradición mágica, una ``ciencia''
(y que tiene que ser puesta entre comillas si
es que pensamos que esa palabra significa lo
que para Platón) oscura a la luz del logos,
acaso en la religión o en las tradiciones
culturales. Con esta tradición se hace
evidente la dicotomía del fármacon: ¿por qué
no se ha extinguido? No hay que ser tan
obstinados como para pensar que todas
aquellas personas se hallan en la locura, y
hay que ser tan humildes como para reconocer
que existen asuntos oscuros a la verdad, que
no se pueden desenredar por completo, y
aprender a vivir con ellos.

Entonces, mi pregunta era si esta reflexión
se tiene desde antes de Derrida, y de qué
manera. No he leído aún la carta para el
amigo japonés, pero justo creo que desde
tiempos antiquísimos, en el oriente se tenían
``filosofías'' cuyas enseñanzas estaban dadas
paradójicamente, dando a entender justamente
ese salto en el que se renuncia a la
explicación y se apela a la vida en un plano
pre-logos.

Creo que Heráclito es el primero que he
estudiado que planteaba explícitamente la paz
con la paradoja, pero más tarde y sobre todo
con la llegada de Aristóteles, aquella idea
se volvió, si acaso ya lo era bastante por su
primer ponente, más oscura que nunca. ¿Qué
pasó en la edad media con todo esto? (ésta sí
es pregunta)

Como sea, lo que se puede ver de todo esto es
que yo creo que es una labor importante de la
filosofía abrir el espacio para una reflexión
que dé más peso a la vida humana que al logos
o a la filosofía misma. Pienso que la
filosofía griega, al menos en
Sócrates-Platón, sí tiene ese acento, pues
parece que la filosofía se dedica en gran
medida procurar el bienestar del hombre.
Desde la misma definición del concepto de
Bien en la República, que se nos figura un
dios filosófico, hasta el problema de la
muerte de Sócrates en el que más vale morir y
hacer el bien que vivir y hacer el mal. Y
esto es en un ámbito tanto social como
personal, o, mejor aún, en el alma.

O sea que yo creo que es importante
desapegarse de la verdad, y por lo pronto
todo esto se asomó del texto de Derrida.

\section{Derrida, Heidegger y el juego}
\label{section-derrida-juego}

The last draft, of June 2018, of my final
essay for Filosofía del Lenguaje. It had no
title yet. One sentence of the draft is
left unfinished.

El ensayo se basa en la crítica de Derrida al
lenguaje occidental. Se trata de mostrar: (1)
mi lectura o manera en la que entiendo la
construcción de Derrida para criticar el
lenguaje; (2) los problemas que yo encontré
en esta tesis, a saber, (2.1) que es una
crítica hecha de lenguaje que critica al
lenguaje y esto quiere decir que se invalida
a sí misma, y (2.2) que parece señalar la
imposibilidad de la comunicación en general;
(3) la resolución de cada uno de estos
problemas: (3.1) a través de la supremacía de
la experiencia sobre el contenido semántico,
y (3.2) a través del ejemplo del pensamiento
oriental; y finalmente (4) una reflexión en
la que se plantea la necesidad de juego o
desplazamiento entre distintos planos de
comunicación y experiencia.

Comenzaré con Derrida, haciendo un repaso de
la construcción que hace para exhibir la
inconsistencia del lenguaje occidental.
Comencemos con la siguiente cita: ``la
escritura, como suplemento del
habla''\footnote{La farmacia de Platón}.
Derrida encuentra a la escritura en una
situación de suplementariedad por lo menos en
la mitología egipcia y en Rousseau. Pero no
nos interesa tanto la explicación de en qué
sentido es la escritura suplemento del habla
como el hecho de que la escritura pueda ser
suplemento del habla: ser suplemento del
habla exhibe la posibilidad de
suplementariedad.

La implicación principal de esto es la
apertura de un juego, por lo pronto en el
sentido de que se puede ganar o perder.
Entendemos el sitio del suplemento como
subordinado, como versión distorsionada,
mentirosa, oscura, falsa, luna, Zot, Hermes.
Derrida agrupa todo eso en una misma figura,
que, en el juego, quiere decir perder.
Derivar la inconsistencia del lenguaje quiere
decir mostrar que la pretensión comunicativa
que tiene el lenguaje nunca se satisface por
completo. Esto es, ningún discurso es
ganador.

El trabajo de Derrida consiste en mostrar que
el suplemento vuelve sobre lo suplementado de
tal forma que ambos nos parezcan
indistinguibles: ``La secundariedad que se
creía poder reservar a la escritura afecta a
todo significado en general, lo afecta desde
siempre, vale decir desde la apertura del
juego''\footnote{De la gramatología, p. 12}.
En el caso de la escritura, terminamos por
ver que el habla realmente no se encuentra
exenta de todos los problemas hermenéuticos
que se plantearon a través de su suplemento.
En sentido amplio, se trata de ver que la
frontera entre la figura de perder y la de
ganar no está bien definida, que el lenguaje
no escapa al ``juego de referencias
significantes''\footnote{De la gramatología
p. 12}, que el juego no concluye.

Vale aquí reafirmar la existencia del
suplemento. Ya hemos dicho que Derrida
encontró a la escritura como suplemento del
habla en distintas instancias, pero, ya
extendida la figura de la escritura, hay que
ver ahora en el pensamiento occidental que el
suplemento existe en el sentido de que los
conceptos a menudo se definen respecto de lo
que no son, por oposición a otros conceptos.
Parte del trabajo de Derrida, y también de
Heidegger, consiste en mostrar que las
lenguas occidentales construyen conceptos de
este modo y que eso genera inconsistencias o
incapacidades lingüísticas.

Esto basta como presentación de la tesis que
quiero analizar. Debemos ahora destacar un
problema que surge naturalmente a partir de
ella: lo que acabamos de repasar es una
construcción hecha de lenguaje en la que se
critica el lenguaje. Si nos convence esta
construcción y somos bastante honestos, lo
primero que debemos hacer es notar que ella
misma es lenguaje y entrar en shock. Éste no
es un problema trivial: nos corresponde
cuestionar cuál es el valor y el alcance de
una crítica filosófica a la capacidad
comunicativa del lenguaje realizada desde el
lenguaje y la filosofía. Parece ser que la
propia tesis está en riesgo de desmantelarse
a sí misma, o quizás confinarse a un terreno
de obsolecencia. En este ensayo quiero
mostrar que esto no es así. Quiero presentar
el terreno en el que debemos situar la tesis
para encontrar su valor.

Comenzaré por decir que el camino para llegar
a este terreno es semejante al camino que se
tomó para construir la tesis: de hecho, el
modo en el que entraremos en paz con la
inconsistencia inmanente que hemos encontrado
estriba en el estado de shock en el que nos
hallamos gracias a su descubrimiento.

Quiero notar aquí que al mismo tiempo que
justificaremos el `suicidio' de la tesis,
justificaremos también el enrredadísimo
lenguaje que utilizan Derrida y Heidegger en
sus textos. Esto es porque el estado de shock
en el que nos hallamos está relacionado con
el shock, desconcierto, insatisfacción,
enojo, que sentimos cuando leemos un texto
que ofrece frases tan desalentadoras como:
``el modo como el habla en tanto que tal
habla ``esencia'', o sea, perdura, es decir,
permanece recogida en aquello que consiente
el habla en su propiedad en tanto que el
habla a sí misma''\footnote{El camino al
habla, p. 225}.

Recibimos esta clase de frases si bien con
humildad, dejando algo de espacio para que
nuestro desconcierto estribe en una
impotencia de nuestra comprensión, también
con una pregunta algo quejumbrosa: ¿acaso es
absolutamente necesario escribir de esta
forma?

Pero puede ser que la respuesta se haya dado
antes de que surgiera la pregunta. En una
conferencia distinta de la recién citada,
pero incluida en el mismo volumen, El camino
al habla, tenemos por primer párrafo:

\begin{quote}
Las tres conferencias que siguen llevan por
título ``La esencia del habla'' quisieran
llevarnos ante la posibilidad de hacer una
experiencia con el habla. Hacer una
experiencia con algo ---sea una cosa, un ser
humano, un dios--- significa que algo nos
acaece, nos alcanza; que se apodera de
nosotros, que nos tumba y nos transforma.
Cuando hablamos de ``hacer'' una experiencia,
esto no significa precisamente que nosotros
la hagamos acaecer; hacer significa aquí:
sufrir, padecer, tomar lo que nos alcanza
receptivamente, aceptar, en la medida en que
nos sometemos a ello. Algo se hace, adviene,
tiene lugar.\footnote{El camino al habla
p.143}
\end{quote}

La advertencia es suficientemente clara. Es
el anuncio del estado de shock que estamos a
punto de enfrentar. Quizás ese estado de
shock es la finalidad del texto mismo, y
entonces sería absolutamente necesario
escribir de esta forma. Pero de nuevo, seamos
humildes, porque la experiencia con el habla
que persigue Heidegger no es precisamente la
frustración que nos produce la
incomprensibilidad de algunas partes de su
texto. Lo que quiero resaltar aquí es lo
siguiente: se está poniendo en primer término
una referencia a la experiencia humana. Esto
es: la pretensión del texto no es mostrar la
verdad. No es resolver un problema o
responder una pregunta. Esto licencia al
discurso a sostener afirmaciones del
siguiente tipo: ``Pese a todo, con estos
enunciados y cualquiera que sea la amplitud
de sus implicaciones, no logramos más que
hacer la suma \ldots{}''\footnote{El camino
al habla, P. 152}. Destacamos en cursivas lo
que nos importa.

Ya se va desvelando el camino. No hay que
empeñarse en la amplitud de las implicaciones
que tiene este discurso tanto como atender a
la experiencia que tenemos cuando lo leemos o
decimos. En particular, una de las
implicaciones podría ser la invalidación del
mismo discurso.

Fijar nuestra mirada en la experiencia y no
en el contenido semántico es una herramienta
perfecta para defendernos de cualquier
crítica lógica. Ésta la primera agresión que
nos parece inquietante, pues el suicidio del
discurso es una reducción al absurdo: de ser
cierta la tesis, en particular ella misma
sería falsa. Pero la reducción absurdo es una
operación dentro del lenguaje, sujeta a sus
reglas; y nosotros atendemos a la experiencia
del lenguaje, que observa sus reglas desde
fuera.

Hablar sobre el estilo de escritura de
Heidegger y Derrida es fundamental para
comprender su filosofía porque no es
accidental ni cruel. El discurso se pliega
sobre sí mismo de una manera verdaderamente
desconcertante porque busca el desconcierto,
porque requiere el desconcierto: el
desconcierto viene a ser clave para entender
la propuesta filosófica porque es
experiencia; está en un plano distinto al de
la significación del discurso. Esto es, el
discurso tiene una carga semántica que ya fue
repasada, pero también tiene una carga
pragmática que llamamos estado de shock, que
es consecuencia de haber mostrado el suicidio
del discurso y al mismo tiempo, de
simplemente haber leído el texto, así como
está escrito. Terminamos por descubrir que el
juego que presenta Derrida, además de ser una
herramienta necesaria para plantear su tesis,
tiene una utilidad quizá más importante:
divertirnos.

Un ejemplo de esto nos lo da Barbara Cassin:
la filosofía se delimita a través de la
sofística, y en tanto que no-sofística, la
filosofía es equivalente a la sofística y no
logra escapar de ella.\footnote{Barbara
Cassin, El efecto sofístico, introducción}
Exhibir que la filosofía se delimita respecto
a la sofística es exhibir el modo en el que
se acredita la verdad filosófica.

Ahora volvamos un poco a la construcción de
conceptos por oposición de la que hablé en el
repaso inicial. Encontramos en El camino al
habla un capítulo que es un diálogo entre un
japonés y un ``inquiridor'':

\begin{quote}
Japonés: Así parece. Porque desde el
encuentro con el pensamiento europeo se ha
evidenciado una incapacidad del habla
nuestra.

Inquiridor: ¿Cómo así?

J: Le falta la fuerza delimitadora con la que
los objetos pueden ser representados los unos
respecto a los demás en un orden claro, es
decir, en relaciones mutuas de jerarquía y
subordinación.

I: ¿Considera usted seriamente esta
incapacidad como un defecto empobrecedor de
su habla?
\end{quote}

La segunda intervención del japonés habla
precisamente sobre este tema. La construcción
de Derrida es exhibir esta relación de
jerarquía y subordinación en los conceptos
mediante la figura de ``la escritura, como
suplemento del habla''. Tenemos en ambos
autores una crítica fuerte al modo en el que
se construyen los conceptos en las lenguas
occidentales y se nos incita, en efecto, a
considerarlo como un defecto empobrecedor de
nuestra habla.

Renunciar a este esquema de construcción de
conceptos es renunciar a la renuncia porque
la renuncia se plantea en términos de esta
construcción conceptual, pero, sobre todo, es
abrir la posibilidad de situarnos en un plano
independiente de esta construcción.

Ya hemos visto cómo el estado de shock
muestra un camino para romper o renunciar a
este esquema: la experiencia humana excede
los límites del campo de la verdad
lingüística. Pero ahora tenemos una
herramienta más, que surge amablemente con la
personalidad japonesa: el pensamiento
oriental, que abre la posibilidad de una
reflexión distinta de lo que Heidegger llama
en ese texto ``el discurso de la
metafísica''\footnote{Camino al habla, p. 96}
--- entendamos por esto la construcción por
oposición de los conceptos.\footnote{Vale
aquí decir que el término `discurso de la
metafísica' es más amplio que la mera
construcción por oposición de la que estamos
hablando. Sin embargo, para que no escape de
nuestras manos la discusión en el terreno
metafísico, nos basta decir que la metafísica
gira en torno del concepto de ser y su
posibilidad significativa, del cual se
desprende la posibilidad de no ser. Aquí
encontramos la construcción por oposición.}

En el apartado anterior hablé del modo en que
nos defendemos de la reducción al absurdo en
términos lógicos. Ahora quiero mostrar otro
problema que surge

Un argumento de la misma índole pero dado en
términos más amplios (y por ello más
peligroso) es que este discurso desemboca en
un terreno relativista en el que el lenguaje
se nos antoja completamente inútil o
impotente en su función comunicativa. Hay que
argumentar que ésta no es la intención de la
tesis.

Renunciar al discurso de la metafísica no
debe entenderse como una resignación o una
derrota, o como la aceptación de que estamos
confinados a un vacío de no-comunicación y de
que debamos desechar la filosofía y el
pensamiento occidental. Al contrario, nos
dice el japonés: ``Para nosotros, el vacío es
el nombre eminente para lo que usted quisiera
decir con la palabra
``ser''''\footnote{Camino al habla, p. 98}.

Esta cita es excelente por dos razones. La
primera ya la mencioné arriba: el japonés
(aunque estuviera callado) entra en escena
como el representante de una clase de humanos
que piensan fuera del esquema metafísico.
Esto es, la personalidad japonesa traza
automáticamente la distinción entre el
pensamiento occidental y oriental.
(Apresurémonos a decir que resaltar esta
distinción nos sitúa en el lado occidental
(oriental=no occidental), más aún, que
nuestra intención es deconstruirlo desde él
mismo; y sigamos adelante). Ver en el
pensamiento oriental una posibilidad racional
distinta de la occidental es importante
porque nadie se atrevería a decir que dicho
pensamiento consta de reflexiones inútiles, o
vacías.

Y si nos atreviéramos, encontraríamos una
interesante sorpresa al releer la cita:
``Para nosotros, el vacío es el nombre
eminente para lo que usted quisiera decir con
la palabra ``ser''''. De hecho, estas
palabras muestran un rasgo del pensamiento
oriental que podemos encontrar en textos muy
antiguos, y es que muchas ideas que a
nosotros nos generan grandes incomodidades
lógicas, como la vacuidad y el no-ser, son
abiertamente aprehendidas.

Veámoslo en en el Budismo, donde el modo para
vivir una vida plena estriba precisamente en
el desapego, que tiene relación directa con
la vacuidad, y por cierto despierta sospechas
en nuestro pensamiento lógico: ¿acaso no
habría entonces que desapegarnos del mismo
Budismo? Ésta no parece ser una inquietud:

\begin{quote}
``[\ldots{}] No hay vista, no hay oído, no
hay olfato \\
Ni gusto ni cuerpo ni mente \\
No hay color, no hay sonido \\
Ni olor ni sabor ni tacto \\
Ni nada a lo que la mente pueda aferrarse \\
Ni siquiera el hecho de sentir \\
{[}\ldots{}] \\
Todos los budas del pasado y del presente \\
Los budas del tiempo futuro \\
Usando esta sabiduría Prajña \\
Alcanzan la visión completa y perfecta \\
{[}\ldots{}]''\footnote{Sutra del corazón}
\end{quote}

Se trata aquí de mostrar que el pensamiento
oriental sostiene un discurso contradictorio
desde la antigüedad, y de hecho es el
predilecto para enseñar su doctrina
espiritual. No es fortuita la elección de
este tipo de discurso. Hay que ver en ella la
solución al problema de cómo transmitir lo
esencial de una profunda realización
espiritual que algunos de sus practicantes
alcanzaron. Éste es un problema para nada
trivial, y de hecho exige una gran destreza
lingüística para su resolución ---como el
problema de la discusión sobre el lenguaje.

Vimos primero que también el juego de Derrida
es una figura que apunta hacia un terreno
práctico a través del estado de shock y lo
divertida que pudiera ser la lectura. Pero
hay que ver en el juego también una figura
por completo seria y aburrida. Ya no jugamos;
ahora creemos ampliamente en las
posibilidades comunicativas del lenguaje para
transmitir la propuesta que las palabras de
Derrida cargan consigo, y la sostenemos: el
Sutra del Corazón nos anima a hacerlo.

Esta tesis, en cuanto que la sostenemos, nos
permite conocer al lenguaje. La tesis de
Derrida sí es una crítica al lenguaje
occidental y sí debemos sostenerla dentro del
campo del lenguaje para obtener de ella el
beneficio filosófico que ofrece. Es decir,
lejos de que la tesis desemboque en un
terreno de vacío (occidental) comunicativo,
como temíamos arriba, la tesis nos habla
(puede hacerlo) sobre el lenguaje, y debemos
aprovechar este discurso: la tesis desemboca
en un vacío oriental comunicativo.

El juego, la figura suicida y contradictoria,
es un trickster de mal gusto. Demos un paso
atrás para ver lo que ha pasado en estas
páginas. Primero, decimos que jugamos el
juego con el fin de divertirnos y así
brincamos cómodamente todas las barreras
lógicas que su planteamiento supone. Luego,
dejamos de jugarlo en cuanto nos parece
necesario para extraer de él un jactancioso
`beneficio filosófico' que consiste en una
crítica al lenguaje occidental.

Demos una vuelta más: esas dos operaciones
conforman un juego. No hemos hecho más que
traer a estas páginas la figura del juego no
a través de una cita o un recuerdo sino a
través de la estructura de nuestro discurso.
Lo dijimos antes: ésta es precisamente la
crítica que tratamos de hacer: que el
discurso no escapa al juego.

La posibilidad de jugar es el alfil de
nuestro discurso porque muestra cómo funciona
el lenguaje. Tenemos la capacidad de creer en
el significado de nuestras palabras, de
aferrarnos a la estructura lógica y gramática
que los símbolos guardan entre sí para
construir un universo semántico que existe en
cuanto creemos en él. Pero también tenemos la
capacidad de no creer, de soltar las ataduras
lógicas que tantas veces son sofocantes, de
no construir ese universo de significados y
verlo desde afuera; y reírnos de él.

Éstas son dos caras de la misma moneda que no
pueden ser separadas. Nunca podremos librar
por completo nuestro discurso de la
pretensión de verdad, de esa impresión,
quizás ilusoria, de que lo que decimos de
hecho significa y puede corresponder con la
realidad. Pero nunca debemos ser tan ingenuos
como para olvidar que todo lo que acabamos de
decir podría simplemente no ser así.

\medskip
The draft was sent to the teacher with a
note, of which this is the philosophical
part.

Durante mucho tiempo tuve la idea
de discutir un poco la frase de Heidegger de
``el lenguaje es la casa del ser'', en el
sentido que yo le he dado, a saber, que es a
través de nuestra capacidad de construir ese
``universo semántico'' que estamos
contruyendo todos los entes que nos rodean.
Hay que poner en cursivas `todos', porque, en
efecto, al momento de asegurar la existencia
de cualquier ente, lo hacemos desde el
lenguaje. No podemos asegurar la existencia
de cualquier cosa más allá del lenguaje
porque `asegurar la existencia de' es un acto
de naturaleza lingüística.

(Pero finalmente quedé un poco desalentado de
hablar sobre eso porque me pareció, sobre
todo después de leer a Lafont, que esa frase
y en general la interpretación del lenguaje
de Heidegger tienen un contenido metafísico
muy profundo que no entiendo del todo. Así
que a veces me da la impresión de que esa
interpretación de ``el lenguaje es la casa
del ser'' se escapa de lo que él quería
decir.)

Como sea, me encantaría comunicar en mi
ensayo que este acto de creación, aunque
parece desmoronar nuestra realidad de manera
espeluznante, de ser debe ser recibido como
algo hermoso, de lo que no podemos, pero
sobre todo nunca quisiéramos, deshacernos;
que no debemos tener miedo de pensar que los
entes que nos rodean son creados a través del
lenguaje, sino aprender a jugar con ellos
desde una perspectiva consciente de la
existencia plano pre-lingüístico. Es mi
opinión personal que estos mundos creados son
maravillosos, ¡incluso los que creamos sin
darnos cuenta y tomamos por verdaderos!

Creo que la elección de Derrida de la figura
del juego, principalemente por su naturaleza
contradictoria, es perfecta para comunicar
este mensaje porque su enseñanza es práctica:
a través de la contradicción vivimos el
desplazamiento desde dentro del mundo creado
hacia afuera (y siempre que salimos, estamos
nuevamente en otro mundo creado). De hecho
creo que por esta cualidad es que el discurso
de la filosofía oriental usa tan
frecuentemente una figura contradicotria como
la vacuidad o el Tao de Lao Tse.


\bibliography{my}
\bibliographystyle{amsalpha}

\end{document}
